\documentclass[12pt]{article}

% ---- Fuente Times New Roman ---
\usepackage{mathptmx}

% ---- Paquetes matemáticos ---
\usepackage{amsmath, amssymb, bm, physics}
\usepackage{mathtools}
\usepackage{siunitx}

% mejoras a amsmath (colchetes, etc.)
% unidades físicas: \SI{5}{\nano\meter}

% ---- Formato ---
\usepackage{geometry}
\usepackage{setspace}
\usepackage{microtype}
\usepackage{parskip}


% ---- Colores (solo para bloques de codigo) ---
\usepackage{xcolor}

% ---- Codigo fuente ---
\usepackage{listings}

% ---- Tablas y figuras ---
\usepackage{booktabs}
\usepackage{array}
\usepackage{tabularx}
\usepackage{float}
\usepackage{graphicx}
\usepackage{caption}

\usepackage{subcaption}

% ---- Bibliografía ---
\usepackage{natbib}

% ---- Idioma ---
\usepackage{babel}

% ---- Hyperref al final, SIN color (hidelinks = negro puro, sin subrayado) ---
\usepackage[hidelinks]{hyperref}

% ---- Márgenes e interlineado ---
\geometry{margin=1in}
\onehalfspacing

% ---- Estilo de codigo (sin morado, colores discretos) ---
\definecolor{codegreen}{rgb}{0.0, 0.45, 0.0}
\definecolor{codegray}{rgb}{0.4, 0.4, 0.4}
\definecolor{codeblue}{rgb}{0.0, 0.0, 0.55}
\definecolor{backcolour}{rgb}{0.97, 0.97, 0.97}

\lstdefinestyle{pythonstyle}{
backgroundcolor=\color{backcolour},
commentstyle=\color{codegray}\itshape,
keywordstyle=\color{codeblue}\bfseries,
numberstyle=\tiny\color{codegray},
stringstyle=\color{codegreen},
basicstyle=\ttfamily\footnotesize,
breakatwhitespace=false,
breaklines=true,
captionpos=b,
keepspaces=true,
numbers=left,
numbersep=5pt,
showspaces=false,
showstringspaces=false,
showtabs=false,
tabsize=2,
frame=single,
rulecolor=\color{codegray},
xleftmargin=2em,
framexleftmargin=1.5em
}

\lstset{style=pythonstyle}

\title{\textbf{De la Polarizaci\'on At\'omica a la Bioimpedancia:}\\
\textbf{Un Modelo Multiescala para la Inferencia}\\[0.3em]
\textbf{de Par\'ametros Fisiol\'ogicos}\\[0.5em]
\large Universidad Aut\'onoma de Nuevo Le\'on}

\author{Juan Manuel Serena Cece\~{n}a}

\date{}

%=================================================================

\begin{document}

\maketitle

\begin{abstract}
La bioimpedancia el\'ectrica constituye una herramienta ampliamente utilizada
para la caracterizaci\'on de sistemas biol\'ogicos; sin embargo, su interpretaci\'on
suele basarse en modelos emp\'iricos que carecen de una conexi\'on directa con los
procesos f\'isicos y fisiol\'ogicos subyacentes. En este trabajo se propone un
modelo multiescala que establece un v\'inculo consistente entre la polarizaci\'on
microsc\'opica de la materia y la respuesta macrosc\'opica de bioimpedancia.

El desarrollo del modelo inicia con la descripci\'on de la interacci\'on entre un
campo el\'ectrico y la materia a nivel at\'omico, mediante el oscilador de Lorentz
y la polarizabilidad dependiente de la frecuencia. A partir de esta base, se
construye la respuesta diel\'ectrica macrosc\'opica utilizando la relaci\'on de
Clausius--Mossotti, incorporando posteriormente efectos din\'amicos mediante
los modelos de relajaci\'on de Debye y Cole--Cole, los cuales permiten capturar
la dispersi\'on observada en sistemas complejos.

El modelo se extiende hacia sistemas biol\'ogicos mediante la incorporaci\'on de
membranas celulares, tratadas como estructuras diel\'ectricas heterog\'eneas, as\'i
como fen\'omenos de polarizaci\'on interfacial descritos por el mecanismo de
Maxwell--Wagner. Asimismo, se consideran efectos fisiol\'ogicos asociados al
transporte i\'onico, particularmente la din\'amica de canales de sodio y potasio,
introduciendo una dependencia temporal y no lineal en la conductividad del sistema.

Desde esta perspectiva, la impedancia compleja se interpreta como un observable
f\'isico que encapsula informaci\'on sobre la din\'amica interna del sistema biol\'ogico.
Con el fin de abordar el problema inverso, se implementa un enfoque basado en

aprendizaje autom\'atico que permite inferir par\'ametros fisiol\'ogicos a partir de
la respuesta en frecuencia generada por el modelo.

Los resultados demuestran que es posible reconstruir propiedades
internas del sistema a partir de mediciones externas de
bioimpedancia. El modelo inverso, basado en una red neuronal
entrenada con datos sintéticos físicamente fundamentados, logra
estimar parámetros fisiológicos clave ---incluyendo la conductividad
de base $\sigma_{\text{base}}$, la conductancia de sodio $g_{Na}$
y el tiempo de relajación $\tau$--- con errores relativos (MAPE)
de \textbf{2.9\%}, \textbf{5.8\%} y \textbf{14.0\%},
respectivamente, sobre el conjunto de prueba. Parámetros
adicionales como la conductancia de potasio $g_K$ y el exponente
Cole-Cole $\alpha_{\text{CC}}$ presentan errores medianos de
9.9\% y 14.0\%, reflejando la mayor dificultad intrínseca de su
identificación a partir del espectro de impedancia. Este trabajo
establece un marco teórico que integra física fundamental,
modelación biológica y herramientas computacionales, proponiendo
una nueva interpretación de la bioimpedancia no solo como
herramienta de medición, sino como una ventana hacia la dinámica
interna de sistemas vivos.
\end{abstract}
\section*{Agradecimientos}
Deseo expresar mi más profundo agradecimiento a todas aquellas personas que, de manera directa o indirecta, contribuyeron al desarrollo de este trabajo y al proceso académico y personal que lo hizo posible.
A Dios, por brindarme fortaleza, guía y perseverancia a lo largo de este camino, y por permitirme culminar una etapa fundamental de mi formación.

A mis padres, Maricela y Manuel, por su apoyo incondicional, sus enseñanzas y el esfuerzo constante que han dedicado a mi crecimiento personal y profesional. Su ejemplo de compromiso, responsabilidad y dedicación ha sido el principal fundamento de mis metas y aspiraciones.


A mi tutor y asesor, el Dr. Eduardo Pérez Tijerina, por su orientación, experiencia y acompañamiento académico. Su apoyo y conocimiento fueron esenciales para la consolidación y desarrollo de esta investigación.

A mis hermanas, Joyce y Jane, y a mi perrito Jackie, por su cariño, motivación y presencia constante, que representaron un impulso importante durante este proceso.

A mi tía Juanis, a mis abuelitas y a mi primo Jony, por sus consejos, apoyo y enseñanzas, recordándome siempre el valor del esfuerzo, la constancia y la dedicación.

A Liliana Cortez, quien fue fundamental en este trabajo de inicio a fin. Me enseñaste que no se puede saber todo en esta vida, y que no es impreciso  no tener siempre las respuestas y que, a veces, es necesario caer antes de aprender a levantarse. Gracias por tu apoyo, paciencia y por acompañarme tanto en los buenos como en los malos momentos.

A los licenciados Irving Rocha, Karla Vázquez, Yamileth Páez, Luis Villarreal y David Palafox, por su amistad, tiempo y ejemplo profesional, los cuales enriquecieron mi formación durante mi estancia en la facultad.

A mis roomies, Daniel y Patricio, por su compañerismo, intercambio de ideas y convivencia, que hicieron más llevadero este proceso académico.

Finalmente, a mi familia y amigos, por su apoyo, confianza y cariño incondicional, así como por acompañarme y motivarme a lo largo de este camino.

\section*{Lista de S\'imbolos}
\addcontentsline{toc}{section}{Lista de S\'imbolos}
\begin{center}
\renewcommand{\arraystretch}{1.4}
\begin{tabular}{@{}lll@{}}
\toprule
\textbf{S\'imbolo} & \textbf{Definici\'on} & \textbf{Unidades} \\
\midrule
$Z(\omega)$ & Impedancia el\'ectrica compleja & $\Omega$ \\
$\varepsilon^*(\omega) = \varepsilon' - j\varepsilon''$ & Permitividad compleja & F/m \\
$\varepsilon'$ & Parte real de la permitividad (almacenamiento) & F/m \\
$\varepsilon''$ & Parte imaginaria (p\'erdidas + conducci\'on) & F/m \\
$\varepsilon_s$ & Permitividad est\'atica ($\omega\to 0$) & F/m \\
$\varepsilon_\infty$ & Permitividad de alta frecuencia & F/m \\
$\sigma$ & Conductividad el\'ectrica & S/m \\
$\omega = 2\pi f$ & Frecuencia angular & rad/s \\
$\tau$ & Tiempo de relajaci\'on diel\'ectrica & s \\

$\alpha_{CC}$ & Par\'ametro de ensanchamiento Cole--Cole & --- \\
$\alpha(\omega)$ & Polarizabilidad at\'omica & C$\cdot$m$^2$/V \\
$\mathbf{p}$ & Momento dipolar el\'ectrico & C$\cdot$m \\
$\mathbf{P}$ & Densidad de polarizaci\'on & C/m$^2$ \\
$C_m$ & Capacitancia de membrana por unidad de \'area & F/m$^2$ \\
$V_m$ & Potencial de membrana & V \\
$E_i$ & Potencial de Nernst del ion $i$ & V \\
$g_i$ & Conductancia del canal i\'onico $i$ & S \\
$d$ & Grosor de la membrana & m \\
$a$ & Radio celular & m \\
$\phi$ & Fracci\'on volum\'etrica celular (hematocrito) & --- \\
$\tau_{MW}$ & Tiempo de relajaci\'on de Maxwell--Wagner & s \\
$k_B$ & Constante de Boltzmann & J/K \\
$R$ & Constante universal de los gases & J/(mol$\cdot$K) \\
$F$ & Constante de Faraday & C/mol \\
$\varepsilon_0$ & Permitividad del vac\'io & F/m \\
\bottomrule
\end{tabular}
\end{center}
\newpage

\tableofcontents

\newpage

%=================================================================

\section{Introducci\'on}

%=================================================================

El estudio de las propiedades el\'ectricas de sistemas biol\'ogicos ha cobrado una
importancia creciente en \'areas como la biof\'isica, la ingenier\'ia biom\'edica y la
medicina diagn\'ostica. En particular, la bioimpedancia el\'ectrica se ha consolidado
como una herramienta no invasiva capaz de caracterizar tejidos a partir de su
respuesta a campos el\'ectricos alternos. Sin embargo, en muchos enfoques
tradicionales, la interpretaci\'on de esta magnitud se basa en modelos emp\'iricos
o circuitales que, si bien son \'utiles desde el punto de vista pr\'actico, carecen
de una conexi\'on directa con los procesos f\'isicos y fisiol\'ogicos subyacentes.

La bioimpedancia es una medida de la oposici\'on que presentan los tejidos
biol\'ogicos al paso de una corriente el\'ectrica alterna. Esta magnitud encapsula
tanto efectos resistivos como capacitivos, los cuales emergen de la estructura
y din\'amica del sistema biolog\'ico.

Esta limitaci\'on motiva la necesidad de desarrollar un marco te\'orico que permita
establecer un v\'inculo consistente entre la f\'isica fundamental y la respuesta
macrosc\'opica medida experimentalmente. En este contexto, surge la pregunta
central que gu\'ia el presente trabajo: \emph{¿es posible inferir propiedades internas
de sistemas biol\'ogicos a partir de mediciones de bioimpedancia, mediante un
modelo f\'isico multiescala que integre desde la polarizaci\'on at\'omica hasta la
din\'amica celular?}

Para abordar esta problem\'atica, en esta tesis se propone un modelo que conecta
diferentes niveles de descripci\'on f\'isica. En el nivel microsc\'opico, la interacci\'on
entre la materia y un campo el\'ectrico se describe a trav\'es de la polarizaci\'on
at\'omica, modelada mediante el oscilador de Lorentz y su correspondiente
polarizabilidad dependiente de la frecuencia. Esta descripci\'on permite entender

c\'omo las cargas ligadas responden din\'amicamente a excitaciones externas,
introduciendo efectos de resonancia y disipaci\'on.

A partir de esta base, se establece el puente hacia la escala macrosc\'opica
mediante relaciones constitutivas como la ecuaci\'on de Clausius--Mossotti,
la cual conecta la polarizabilidad microsc\'opica con la permitividad efectiva
del medio. Este paso es fundamental, ya que revela que la respuesta diel\'ectrica
de un material no es simplemente la suma de contribuciones individuales,
sino el resultado de interacciones colectivas entre dipolos.

Posteriormente, se incorporan efectos din\'amicos a trav\'es de modelos de
relajaci\'on como el modelo de Debye y su generalizaci\'on de Cole--Cole,
los cuales permiten describir la dependencia temporal y frecuencial de la
polarizaci\'on en sistemas complejos. En particular, estos modelos capturan
la presencia de m\'ultiples escalas temporales asociadas a procesos f\'isicos
como la rotaci\'on dipolar, la reorganizaci\'on estructural y la heterogeneidad
del medio.

El marco te\'orico se extiende hacia sistemas biol\'ogicos mediante la
incorporaci\'on de la estructura y propiedades de la membrana celular.
Las membranas, compuestas principalmente por bicapas lip\'idicas,
act\'uan como barreras diel\'ectricas altamente selectivas que separan
regiones con distintas propiedades el\'ectricas. Esta configuraci\'on
da lugar a fen\'omenos de polarizaci\'on interfacial, descritos mediante
el mecanismo de Maxwell--Wagner, los cuales son responsables de
caracter\'isticas distintivas en la respuesta de bioimpedancia.

Adicionalmente, la presencia de canales i\'onicos, particularmente de
sodio y potasio, introduce una dependencia din\'amica en la conductividad

del sistema. La apertura y cierre de estos canales, regulados por el
potencial de membrana, generan un comportamiento no lineal que influye
directamente en la respuesta el\'ectrica del tejido. De esta manera, la
bioimpedancia no solo refleja propiedades estructurales, sino tambi\'en
procesos fisiol\'ogicos activos.

Desde esta perspectiva, la impedancia compleja $Z(\omega)$ puede ser
interpretada como un observable f\'isico que encapsula informaci\'on sobre
la din\'amica interna del sistema. Cada uno de sus componentes contiene
contribuciones de distintos mecanismos, incluyendo conducci\'on i\'onica,
polarizaci\'on dipolar y acumulaci\'on de carga en interfaces.

Sin embargo, la relaci\'on entre la respuesta en frecuencia y los par\'ametros
internos del sistema es altamente no lineal y, en general, no invertible de
forma anal\'itica. Esto da lugar a un problema inverso intr\'insecamente complejo,
en el cual m\'ultiples configuraciones f\'isicas pueden producir respuestas
similares. Para abordar esta dificultad, se introduce el uso de t\'ecnicas de
aprendizaje autom\'atico, las cuales permiten aproximar dicha relaci\'on a partir
de datos generados mediante el modelo f\'isico.

El enfoque propuesto en este trabajo se distingue por integrar de manera
consistente la f\'isica fundamental, la modelaci\'on biol\'ogica y herramientas
computacionales modernas, dando lugar a un marco multiescala capaz de
relacionar directamente la bioimpedancia con par\'ametros fisiol\'ogicos.

En conjunto, esta tesis establece un puente entre la descripci\'on microsc\'opica
de la materia y la respuesta macrosc\'opica de sistemas biol\'ogicos, proponiendo
una nueva forma de interpretar la bioimpedancia no solo como una herramienta
de medici\'on, sino como un medio para acceder a la din\'amica interna de los

sistemas vivos.

% ---- Frase central ---
\begin{center}
\textit{
Este trabajo demuestra que la bioimpedancia no es \'unicamente una cantidad medible,
sino un observable que codifica el estado fisiol\'ogico interno de las c\'elulas.
}
\end{center}

% ---- Originalidad ---
\begin{center}
\textit{
Hasta donde es de nuestro conocimiento, este es uno de los primeros intentos de
acoplar completamente modelos electrodin\'amicos, fisiolog\'ia de membrana y aprendizaje
autom\'atico en un \'unico marco predictivo y autoconsistente.
}
\end{center}

\subsection*{Hip\'otesis Central}

\emph{La hip\'otesis central de este trabajo es que las mediciones de bioimpedancia
pueden vincularse de forma directa y rigurosa con par\'ametros f\'isicos microsc\'opicos
---grosor de la membrana, conductividad del citoplasma y fracci\'on volum\'etrica
celular--- a trav\'es de un modelo electrodin\'amico multiescala, y que este v\'inculo
puede invertirse mediante aprendizaje autom\'atico para recuperar par\'ametros
fisiol\'ogicos a partir de una medici\'on el\'ectrica sencilla.}

\subsection*{Alcance del Marco}

El objetivo de este trabajo es desarrollar un modelo multiescala capaz de
conectar la polarizaci\'on at\'omica con la bioimpedancia macrosc\'opica, permitiendo
la inferencia de par\'ametros fisiol\'ogicos a partir de mediciones el\'ectricas.
Este trabajo no se limita al ajuste de datos experimentales, sino que proporciona
un marco f\'isicamente fundamentado donde cada par\'ametro corresponde a una magnitud
medible, y el modelo puede predecir c\'omo cambia el espectro de impedancia ante
variaciones fisiol\'ogicas no observadas durante el entrenamiento.

La relevancia de esta construcci\'on radica en que cada nivel de descripci\'on no es
independiente, sino que constituye una manifestaci\'on emergente del nivel anterior.
De este modo, la bioimpedancia macrosc\'opica puede entenderse como una consecuencia
directa de procesos microsc\'opicos de polarizaci\'on, convirtiendo este trabajo en
genuinamente multiescala y no solo encadenado.

\subsection*{Objetivos}

Este trabajo persigue cuatro objetivos espec\'ificos:

\begin{enumerate}
\item Derivar, desde primeros principios, una cadena de modelos f\'isicamente
consistente que conecte la polarizabilidad at\'omica con la bioimpedancia macrosc\'opica.
\item Implementar esta cadena num\'ericamente y generar una base de datos de espectros
de impedancia con fundamento f\'isico.
\item Entrenar y validar modelos de aprendizaje autom\'atico que resuelvan el problema
inverso de bioimpedancia, habilitando la estimaci\'on de par\'ametros a partir de
datos de impedancia.
\item Extender el marco hacia una descripci\'on fisiol\'ogica que incorpore el potencial
de membrana, gradientes i\'onicos y din\'amica de canales, estableciendo un v\'inculo

conceptual entre f\'isica y biolog\'ia celular.
\end{enumerate}

\subsection*{Estructura del Trabajo}

El trabajo sigue la jerarqu\'ia f\'isica del problema.
La secci\'on~\ref{sec:background} revisa el estado del arte.
La secci\'on~\ref{sec:methods} describe la metodolog\'ia y las definiciones de par\'ametros.
Las secciones~\ref{sec:atomic}--\ref{sec:macro} desarrollan el marco te\'orico.
La secci\'on~\ref{sec:numerical} presenta la implementaci\'on num\'erica y la integraci\'on
de aprendizaje autom\'atico.
La secci\'on~\ref{sec:bio_foundations} introduce las bases biol\'ogicas y electroqu\'imicas
de la membrana.
La secci\'on~\ref{sec:physiological} desarrolla la extensi\'on fisiol\'ogica y bioqu\'imica.
La secci\'on~\ref{sec:results} discute los resultados e interpretaci\'on f\'isica.
La secci\'on~\ref{sec:limitations} aborda las limitaciones del modelo, y la
secci\'on~\ref{sec:conclusions} resume las contribuciones principales.

\section{Antecedentes y Estado del Arte}

\label{sec:background}

%=================================================================

\subsection{Desarrollo Hist\'orico de la Bioimpedancia}

El estudio de las propiedades el\'ectricas de los tejidos biol\'ogicos tiene una historia de m\'as
de un siglo. Las primeras mediciones de Fricke (1925) demostraron que las suspensiones de
eritrocitos se comportan como elementos capacitivos a bajas frecuencias, proporcionando la

primera evidencia experimental del comportamiento diel\'ectrico controlado por la membrana

\citep{Fricke1925, Schwan1957}. Estas observaciones fueron sistematizadas por Herman P.\
Schwan, quien introdujo el concepto de dispersiones diel\'ectricas en materiales biol\'ogicos
y defini\'o las ahora est\'andar regiones de dispersi\'on $\alpha$, $\beta$ y $\gamma$

\citep{Schwan1957, Schwan1994}. Cole y Cole (1940) introdujeron una generalizaci\'on emp\'irica
del modelo de relajaci\'on de Debye para describir los espectros de dispersi\'on ensanchados de
materiales biol\'ogicos \citep{Cole1940}, que sigue siendo la herramienta de ajuste est\'andar
hasta hoy.

\subsection{Dispersiones Diel\'ectricas en Tejidos Biol\'ogicos}

\label{sec:dispersiones}

Los tejidos biol\'ogicos exhiben m\'ultiples regiones de frecuencia distintas en las que sus
propiedades diel\'ectricas cambian significativamente. Estas se denominan \emph{regiones de
dispersi\'on} y se etiquetan $\alpha$, $\beta$ y $\gamma$ en orden creciente de frecuencia
\citep{Schwan1957, Foster1989}.

\subsubsection{Dispersi\'on \texorpdfstring{$\alpha$}{Alfa} (Hz -- kHz)} La dispersi\'on $\alpha$
ocurre a las frecuencias m\'as bajas ($1\,\text{Hz}$--$10\,\text{kHz}$).

Se origina en la relajaci\'on de capas de difusi\'on i\'onica alrededor de superficies membranales
cargadas y en la polarizaci\'on de la atm\'osfera de contraiones que rodea a las macromol\'eculas.

Los tiempos de relajaci\'on asociados son largos ($\tau \sim 10^{-3}$--$10^{-1}\,\text{s}$),
reflejando una lenta reorganizaci\'on i\'onica.

\subsubsection{Dispersi\'on \texorpdfstring{$\beta$}{Beta} (kHz -- MHz)} La dispersi\'on $\beta$
es la m\'as relevante para la espectroscop\'ia de bioimpedancia de c\'elulas y sangre. Ocurre entre
$10\,\text{kHz}$ y $100\,\text{MHz}$ y est\'a gobernada por el
mecanismo
de polarizaci\'on interfacial de Maxwell--Wagner \citep{Maxwell1873, Wagner1914}. La corriente
transita de fluir alrededor de las c\'elulas a bajas frecuencias a penetrar el interior celular
a frecuencias m\'as altas. El tiempo de relajaci\'on caracter\'istico es:

\begin{equation}
\tau_{MW} = \frac{\varepsilon_1 + 2\varepsilon_3}{\sigma_1 + 2\sigma_3}
\sim 10^{-7}\text{--}10^{-5}\,\text{s}
\end{equation}

Esta dispersi\'on codifica la capacitancia de la membrana, la conductividad del citoplasma y
la geometr\'ia celular, convirti\'endola en el objetivo primario de este trabajo.

\subsubsection{Dispersi\'on \texorpdfstring{$\gamma$}{Gamma} (GHz)} La dispersi\'on
$\gamma$ ocurre por encima de $1\,\text{GHz}$ y est\'a dominada por la relajaci\'on rotacional
de las mol\'eculas de agua ($\tau \approx 8\,\text{ps}$ a
$25\,^\circ\text{C}$)

\citep{Kaatze1989}.

\begin{table}[H]
\centering
\caption{Resumen de las regiones de dispersi\'on diel\'ectrica en tejidos biol\'ogicos.}
\label{tab:dispersiones}
\begin{tabular}{@{}llll@{}}
\toprule
\textbf{Dispersi\'on} & \textbf{Rango de frecuencia}
& \textbf{Mecanismo principal} & \textbf{Tiempo de relajaci\'on} \\
\midrule
$\alpha$ & $1\,\text{Hz}$--$10\,\text{kHz}$
& Difusi\'on i\'onica, polarizaci\'on de contraiones
& $10^{-3}$--$10^{-1}\,\text{s}$ \\
$\beta$ & $10\,\text{kHz}$--$100\,\text{MHz}$
& Polarizaci\'on interfacial Maxwell--Wagner
& $10^{-7}$--$10^{-5}\,\text{s}$ \\
$\gamma$ & $>1\,\text{GHz}$
& Rotaci\'on del dipolo del agua
& $\sim 8\,\text{ps}$ \\
\bottomrule
\end{tabular}
\end{table}

\subsection{Modelos Existentes y sus Limitaciones}

Varios modelos de circuito equivalente han sido propuestos para describir la respuesta

el\'ectrica de tejidos biol\'ogicos. El modelo de impedancia de Cole \citep{Cole1940} es el
m\'as com\'un:

\begin{equation}
Z(\omega) = R_\infty + \frac{R_0 - R_\infty}{1 + (j\omega\tau)^\beta}
\end{equation}

Si bien proporciona excelentes ajustes a datos experimentales, es puramente fenomenol\'ogico:
los par\'ametros $R_0$, $R_\infty$ y $\tau$ son coeficientes de ajuste sin significado fisiol\'ogico
directo. No pueden predecir c\'omo cambia $Z(\omega)$ cuando, por ejemplo, el
hematocrito disminuye en la anemia, ni pueden generar datos sint\'eticos f\'isicamente realistas
para aprendizaje autom\'atico. Esto motiva el enfoque basado en f\'isica del presente trabajo.

%=================================================================

\section{Metodolog\'ia}

\label{sec:methods}

%=================================================================

\subsection{Estrategia General de Modelado}

Este trabajo sigue una jerarqu\'ia f\'isica ascendente (\textit{bottom-up}). Cada nivel de
descripci\'on se deriva del anterior, garantizando consistencia interna entre escalas:

\begin{center}

\renewcommand{\arraystretch}{1.3}
\begin{tabular}{@{}ccccccc@{}}
Cu\'antica & $\to$ & Cl\'asica & $\to$ & Mesosc\'opica & $\to$ & Macrosc\'opica \\
$\alpha_{\text{QM}}$ & & $\varepsilon^*(\omega)$ & &
$Z_{\text{c\'elula}}$ & & $Z_{\text{sangre}} \to$ ML
\end{tabular}
\end{center}

El aprendizaje autom\'atico no es una herramienta independiente en este marco. Es una
extensi\'on del modelo f\'isico que permite resolver problemas inversos anal\'iticamente
intratables. Sin el
modelo f\'isico proporcionando datos sint\'eticos de entrenamiento con par\'ametros verdaderos
conocidos, el componente de aprendizaje autom\'atico no tendr\'ia una base confiable. Ambos
son inseparables.

\subsection{Nivel de descripci\'on del modelo}
\label{sec:nivel_descripcion}

El presente trabajo se desarrolla dentro de una
\emph{descripci\'on efectiva macrosc\'opica}, en la cual
los efectos microsc\'opicos son incorporados mediante
par\'ametros constitutivos dependientes del estado
fisiol\'ogico. Esta aproximaci\'on permite conectar
fen\'omenos electromagn\'eticos con propiedades biol\'ogicas
sin resolver expl\'icitamente la estructura microsc\'opica.

El modelo articula tres niveles de descripci\'on:
\begin{itemize}
\item \textbf{Microsc\'opico}: polarizabilidad at\'omica
(Lorentz, Clausius--Mossotti).
\item \textbf{Mesosc\'opico}: permitividad efectiva
Cole--Cole, membrana (Pauly--Schwan,
Maxwell--Garnett completo).
\item \textbf{Macrosc\'opico}: $Z(\omega, t)$ acoplada
a la din\'amica de Hodgkin--Huxley.
\end{itemize}
La conexi\'on entre niveles es derivada desde primeros
principios, no fenomenol\'ogica.

\subsection{Estructura Dual del Modelo}

Este trabajo se estructura a lo largo de dos direcciones complementarias:
\begin{itemize}
\item \textbf{Marco f\'isico fundamental}: deriva la din\'amica de la membrana desde primeros
principios, conectando polarizabilidad at\'omica con bioimpedancia macrosc\'opica.
\item \textbf{Enfoque basado en datos}: habilita la predicci\'on de estados biol\'ogicos
a partir de se\~{n}ales el\'ectricas medibles mediante aprendizaje autom\'atico.
\end{itemize}

Ambos enfoques est\'an interconectados y se refuerzan mutuamente.

\subsection{Validez de la Aproximaci\'on Cuasiestática}

A lo largo de este trabajo se emplea la aproximaci\'on cuasiestática, que reduce las ecuaciones
de Maxwell a la ecuaci\'on de Laplace. Esto es v\'alido cuando:

\begin{equation}
\lambda = \frac{c}{f\sqrt{\varepsilon_r}} \gg L_{\text{c\'elula}}
\end{equation}

Para $a \sim 5\,\mu\text{m}$ y $f = 1\,\text{GHz}$, $\varepsilon_r \sim 80$:

\begin{equation}
\lambda_{\min} = \frac{3\times10^8}{10^9\cdot\sqrt{80}} \approx 33\,\mu\text{m}
\gg 5\,\mu\text{m}
\end{equation}

La condici\'on se satisface en m\'as de tres \'ordenes de magnitud a lo largo de todo el rango de
frecuencias. Los efectos de retardo y las p\'erdidas por radiaci\'on se desprecian con seguridad.

\subsection{Definici\'on de S\'imbolos y Espacio de Par\'ametros F\'isicos}

Todos los s\'imbolos utilizados en este trabajo se definen en la Tabla~\ref{tab:simbolos}.

Los par\'ametros f\'isicos con sus rangos fisiol\'ogicos y valores de simulaci\'on se presentan en
la Tabla~\ref{tab:params}.

\begin{table}[H]
\centering
\caption{Definici\'on completa de s\'imbolos empleados a lo largo de este trabajo.}
\label{tab:simbolos}
\renewcommand{\arraystretch}{1.3}
\begin{tabular}{@{}lll@{}}
\toprule
\textbf{S\'imbolo} & \textbf{Definici\'on} & \textbf{Unidades} \\
\midrule
$\varepsilon_0$ & Permitividad del vac\'io ($8.854\times10^{-12}$) & F/m \\
$\varepsilon_r$ & Permitividad relativa & --- \\
$\varepsilon_i$ & Permitividad absoluta de la regi\'on $i$ & F/m \\
$\varepsilon^*$ & Permitividad compleja: $\varepsilon' - j\varepsilon''$ & F/m \\
$\varepsilon'$ & Parte real (almacenamiento de energ\'ia) & F/m \\
$\varepsilon''$ & Parte imaginaria (p\'erdidas + conducci\'on) & F/m \\
$\varepsilon_s$ & Permitividad est\'atica (l\'imite baja $f$) & F/m \\
$\varepsilon_\infty$ & Permitividad de alta frecuencia & F/m \\
$\sigma$ & Conductividad el\'ectrica & S/m \\
$\omega = 2\pi f$ & Frecuencia angular & rad/s \\
$\tau$ & Tiempo de relajaci\'on diel\'ectrica & s \\
$\alpha_{CC}$ & Par\'ametro de ensanchamiento Cole--Cole & --- \\
$e$ & Carga elemental ($1.602\times10^{-19}$) & C \\
$m$ & Masa del electr\'on ($9.109\times10^{-31}$) & kg \\
$\omega_0$ & Frecuencia natural del electr\'on ligado & rad/s \\
$\gamma$ & Coeficiente de amortiguamiento & rad/s \\
$\mathbf{p}$ & Momento dipolar el\'ectrico & C$\cdot$m \\

$\mathbf{P}$ & Densidad de polarizaci\'on macrosc\'opica & C/m$^2$ \\
$\chi_e$ & Susceptibilidad el\'ectrica & --- \\
$\alpha(\omega)$ & Polarizabilidad at\'omica/molecular & C$\cdot$m$^2$/V \\
$Z(\omega)$ & Impedancia compleja & $\Omega$ \\
$\Phi$ & Potencial el\'ectrico & V \\
$C_m$ & Capacitancia de membrana por unidad de \'area & F/m$^2$ \\
$Z_m(\omega)$ & Impedancia superficial de membrana & $\Omega\cdot$m$^2$ \\
$\phi$ & Fracci\'on volum\'etrica celular (hematocrito) & --- \\
$a$ & Radio celular & m \\
\bottomrule
\end{tabular}
\end{table}
\begin{tabular}{@{}lll@{}}
\toprule
\textbf{S\'imbolo} & \textbf{Definici\'on} & \textbf{Unidades} \\
\midrule
$d$ & Grosor de la membrana & m \\
$k_B$ & Constante de Boltzmann ($1.381\times10^{-23}$) & J/K \\
$T$ & Temperatura absoluta & K \\
$D_r$ & Coeficiente de difusi\'on rotacional & rad$^2$/s \\
$V_m$ & Potencial de membrana & V \\
$E_i$ & Potencial de Nernst del ion $i$ & V \\
$g_i$ & Conductancia del canal i\'onico $i$ & S \\
$I_i$ & Corriente i\'onica del ion $i$ & A \\
$R$ & Constante de los gases ideales & J/(mol$\cdot$K) \\
$F$ & Constante de Faraday & C/mol \\
$z_i$ & Valencia i\'onica & --- \\
\bottomrule
\end{tabular}
\

\begin{table}[H]
\centering
\caption{Par\'ametros f\'isicos del modelo celular, rangos fisiol\'ogicos y valores de simulaci\'on utilizados en este trabajo \citep{Schwan1957, Asami2002}.}
\label{tab:params}
\renewcommand{\arraystretch}{1.35}
\begin{tabularx}{\textwidth}{@{}lXcc@{}}
\toprule
\textbf{S\'imbolo} & \textbf{Descripci\'on} & \textbf{Rango fisiol.}
& \textbf{Valor sim.} \\
\midrule
$\varepsilon_1/\varepsilon_0$ & Permitividad relativa del citoplasma
& 50--80 & 60 \\
$\sigma_1$ (S/m) & Conductividad del citoplasma
& 0.3--1.5 & 0.5 \\
$\varepsilon_2/\varepsilon_0$ & Permitividad relativa de la membrana
& 2--10 & 5 \\
$\sigma_2$ (S/m) & Conductividad de la membrana
& $10^{-8}$--$10^{-6}$ & $10^{-7}$ \\
$\varepsilon_3/\varepsilon_0$ & Permitividad relativa del plasma
& 70--90 & 80 \\
$\sigma_3$ (S/m) & Conductividad del plasma
& 0.8--1.8 & 1.2 \\
$d$ (nm) & Grosor de la membrana
& 3--10 & 5 \\
$a$ ($\mu$m) & Radio celular (eritrocito)
& 3--8 & 5 \\
$\phi$ & Hematocrito (fracci\'on volum\'etrica)
& 0.35--0.55 & 0.40 \\
$\tau$ (s) & Tiempo de relajaci\'on Cole--Cole
& $10^{-7}$--$10^{-5}$ & $10^{-6}$ \\
$\alpha_{CC}$ & Par\'ametro de ensanchamiento Cole--Cole

& 0.05--0.25 & 0.10 \\
\bottomrule
\end{tabularx}
\end{table}

\subsection{Estrategia de Aprendizaje Autom\'atico}

El componente de aprendizaje autom\'atico aborda dos tareas estrechamente vinculadas:
\paragraph{Problema directo.} Dado el conjunto de par\'ametros f\'isicos
$(\sigma_1, \sigma_3, \phi, d, a)$, predecir el espectro de impedancia $Z(\omega)$.

\paragraph{Problema inverso.} Dado un espectro de impedancia $Z(\omega)$, recuperar los
par\'ametros f\'isicos subyacentes. Esta es la tarea cl\'inicamente relevante, que permite la
estimaci\'on no invasiva del hematocrito, la conductividad del citoplasma y las propiedades
de la membrana. Dado que el mapeo $Z(\omega) \to \{\sigma_1, \sigma_3, \phi\}$ es no lineal
y anal\'iticamente intratable, se entrena una red neuronal para aproximarlo.

%=================================================================

\section{Fundamentos Microsc\'opicos: Polarizaci\'on At\'omica}

\label{sec:atomic}

%=================================================================

\noindent\textit{Esta secci\'on establece el origen microsc\'opico del comportamiento
diel\'ectrico. Toda magnitud macrosc\'opica introducida posteriormente ---permitividad,
polarizaci\'on, impedancia--- se remonta en \'ultima instancia al proceso aqu\'i descrito: el
desplazamiento de un electr\'on ligado por un campo el\'ectrico externo.}

\subsection{Formaci\'on de un Dipolo At\'omico}

Cons\'iderese un \'atomo neutro con carga nuclear $+Ze$ y nube electr\'onica $-Ze$. En equilibrio,
ambos centros coinciden y no existe dipolo. Bajo un campo externo $\mathbf{E}$, la nube
electr\'onica se desplaza, generando un momento dipolar inducido:

\begin{equation}
\mathbf{p} = q\,\mathbf{d}
\end{equation}

donde $q = -e$ y $\mathbf{d}$ es el desplazamiento entre los centros de carga. La constante de
proporcionalidad entre $\mathbf{p}$ y $\mathbf{E}$ define la \emph{polarizabilidad}
$\alpha$, la magnitud central de esta secci\'on.

\subsection{Modelo Cl\'asico: El Oscilador de Lorentz}

Se modela el electr\'on ligado como un oscilador arm\'onico amortiguado impulsado por el
campo. Esta ecuaci'on describe el movimiento de un electr'on ligado al n'ucleo bajo
la acci'on de un campo el'ectrico externo, y es el punto de partida de toda la

construcci'on multiescala de este trabajo:

\begin{equation}
m\ddot{x} + m\gamma\dot{x} + m\omega_0^2\,x = -eE(t)
\label{eq:lorentz}
\end{equation}

\subsubsection{Significado F\'isico de Cada T\'ermino}

Cada t\'ermino en la Ec.~\eqref{eq:lorentz} tiene un papel f\'isico distinto:

\begin{itemize}
\item $m\ddot{x}$ --- \textbf{inercia}: determina c\'omo acelera el electr\'on. Domina a
altas frecuencias; el electr\'on no puede seguir el campo r\'apidamente oscilante
y la polarizaci\'on tiende a cero.
\item $m\gamma\dot{x}$ --- \textbf{amortiguamiento}: representa la disipaci\'on de energ\'ia
(colisiones, radiaci\'on). Responsable de la parte imaginaria de $\alpha(\omega)$
y por tanto de las p\'erdidas diel\'ectricas.
\item $m\omega_0^2 x$ --- \textbf{fuerza restauradora}: la ligadura del electr\'on al
n\'ucleo. Define la frecuencia natural de resonancia $\omega_0$. Domina a bajas
frecuencias, permitiendo el seguimiento cuasiestático del campo.
\item $-eE(t)$ --- \textbf{fuerza motriz}: el campo externo que desplaza al electr\'on
del equilibrio.
\end{itemize}

Los tres reg\'imenes de frecuencia que emergen de esta competencia son:

\begin{itemize}
\item $\omega \ll \omega_0$: domina la fuerza restauradora; el electr\'on sigue el campo
cuasiestáticamente; $\alpha \approx e^2/(m\omega_0^2)$ --- real y constante.
\item $\omega \approx \omega_0$: resonancia; respuesta m\'axima; fuerte absorci\'on.
\item $\omega \gg \omega_0$: domina la inercia; el electr\'on no puede seguir el campo;
$\alpha \to 0$.
\end{itemize}

En el contexto de la bioimpedancia ($f \leq 1\,\text{GHz}$), las resonancias at\'omicas se
encuentran en el ultravioleta ($\omega_0 \sim 10^{16}\,\text{rad/s}$), de modo que siempre se
opera en el r\'egimen $\omega \ll \omega_0$.

\subsubsection{Soluci\'on mediante Representaci\'on Compleja}

Para $E(t) = E_0 e^{j\omega t}$ y $x(t) = x_0 e^{j\omega t}$, sustituyendo en la ecuaci\'on de
movimiento se obtiene la amplitud de desplazamiento:

\begin{equation}
x_0 = \frac{-eE_0/m}{\omega_0^2 - \omega^2 - j\gamma\omega}
\end{equation}

La expresi\'on anterior describe c\'omo el electr\'on responde al campo externo en funci\'on de la
frecuencia. El denominador complejo introduce tanto la resonancia ($\omega \approx
\omega_0$) como la disipaci\'on ($\gamma$). A partir de $p = -ex_0$, se obtiene directamente la
polarizabilidad at\'omica:

\begin{equation}
\boxed{
\alpha(\omega) = \frac{p}{E_0} = \frac{e^2/m}{\omega_0^2 - \omega^2 - j\gamma\omega}
}
\label{eq:alpha}
\end{equation}

donde $\omega_0$ es la frecuencia natural de la ligadura electr\'on--n\'ucleo, $\gamma$ es el
coeficiente de amortiguamiento que caracteriza las p\'erdidas por colisiones y radiaci\'on, y $e$,
$m$ son la carga y la masa del electr\'on, respectivamente.

La polarizabilidad describe la capacidad del 'atomo para deformar su
distribuci'on de carga en respuesta a un campo el'ectrico. Se observa que la
respuesta depende fuertemente de la frecuencia del campo aplicado. En particular,
cuando $\omega \approx \omega_0$, el sistema entra en resonancia, produciendo
una respuesta amplificada. El t'ermino $j\gamma\omega$ introduce un desfase
entre el campo aplicado y la respuesta del sistema, caracter'istico de medios
disipativos. Esta dependencia en frecuencia es la que, propagada hacia escalas
mayores, originar'a la dispersi'on diel'ectrica observada en la bioimpedancia.

\subsubsection{Interpretaci\'on F\'isica de la Polarizabilidad}

La estructura compleja de la Ec.~\eqref{eq:alpha} codifica dos efectos f\'isicos:

\begin{itemize}
\item $\text{Re}(\alpha)$: respuesta en fase; energ\'ia almacenada en el electr\'on
desplazado; contribuye al \'indice de refracci\'on.
\item $\text{Im}(\alpha)$: respuesta en cuadratura; energ\'ia disipada por ciclo;
contribuye a la absorci\'on y las p\'erdidas diel\'ectricas.
\end{itemize}

La idea clave es que la polarizaci\'on \emph{no es instant\'anea}: el amortiguamiento
$j\gamma\omega$ introduce un desfase entre campo y dipolo. Este es el origen at\'omico de la
relajaci\'on, un tema que reaparecer\'a a cada escala del problema.

\subsection{Aproximaci\'on Hidrogenoide Efectiva en Medios Complejos}

En f\'isica at\'omica, los sistemas complejos se simplifican mapeándolos en modelos
hidrogenoides con par\'ametros efectivos. Un enfoque similar puede adoptarse en sistemas
biol\'ogicos introduciendo unidades dipolares efectivas que representan estructuras moleculares
complejas como cadenas lip\'idicas y prote\'inas. Cada unidad se caracteriza por una
polarizabilidad efectiva $\alpha_{\text{eff}}$, que captura la respuesta electr\'onica y
estructural colectiva del sistema. Esta aproximaci\'on permite extender los modelos
electrodin\'amicos microsc\'opicos a materiales biol\'ogicamente relevantes sin resolver
expl\'icitamente toda la estructura molecular.

\subsection{Polarizabilidad Cu\'antico-Mec\'anica}

El modelo cl\'asico predice correctamente la dependencia en frecuencia, pero trata $\omega_0$
como par\'ametro libre. La mec\'anica cu\'antica lo deriva de los niveles de energ\'ia. Usando la
teor\'ia de perturbaciones de segundo orden con $H' = -\hat{d}\cdot E$:

\begin{equation}
\alpha_{\text{QM}} =
2\sum_{n\neq 0}\frac{|\langle n|\hat{d}|0\rangle|^2}{E_n - E_0}
\end{equation}

donde $|\langle n|\hat{d}|0\rangle|^2$ son los elementos de matriz de dipolo de transici\'on y
$(E_n - E_0)$ son las energ\'ias de excitaci\'on.

%=================================================================

\section{De los \'Atomos a los Medios Macrosc\'opicos}

\label{sec:macro_theory}

%=================================================================

\noindent\textit{Habiendo establecido la polarizabilidad de un \'atomo individual, ahora se
pregunta: ¿cu\'al es la respuesta diel\'ectrica colectiva de $N \sim 10^{28}$ \'atomos por metro
c\'ubico? La respuesta no es simplemente $N$ veces la respuesta at\'omica, porque cada \'atomo
es tambi\'en impulsado por los campos de sus vecinos. Esta secci\'on aborda ese problema de
muchos cuerpos y llega a la permitividad macrosc\'opica.}

\subsection{Densidad de Polarizaci\'on y Susceptibilidad}

Para $N$ \'atomos por unidad de volumen, la polarizaci\'on macrosc\'opica es:

\begin{equation}
\mathbf{P} = N\mathbf{p} = N\alpha\,\mathbf{E}_{\text{local}}
\end{equation}

La expresi\'on anterior vincula la respuesta microsc\'opica de un \'atomo
individual ($\alpha$) con la polarizaci\'on macrosc\'opica del medio ($\mathbf{P}$).
Comparando con la definici\'on est\'andar de susceptibilidad, se identifican:

\begin{equation}
\mathbf{P} = \varepsilon_0\chi_e\,\mathbf{E}, \qquad
\chi_e = \frac{N\alpha}{\varepsilon_0}, \qquad
\varepsilon_r = 1 + \chi_e
\end{equation}

La sutileza clave es que $\mathbf{E}_{\text{local}}$ ---el campo que realmente impulsa a cada
\'atomo--- difiere del campo aplicado macrosc\'opico $\mathbf{E}$ debido a las
contribuciones de los dipolos vecinos.

\subsection{Correcci\'on del Campo Local de Lorentz}

Lorentz demostr\'o que el campo local en una mol\'ecula dentro de un medio polarizado es:

\begin{equation}
\mathbf{E}_{\text{local}} = \mathbf{E} + \frac{\mathbf{P}}{3\varepsilon_0}
\end{equation}

El campo local representa el campo efectivo experimentado por una mol\'ecula
dentro del medio. Este no coincide con el campo externo aplicado, debido a
la contribuci'on de los dipolos circundantes. El segundo t'ermino,
$P/(3\varepsilon_0)$, es el \emph{campo de cavidad}: la contribuci\'on de
todos los dipolos circundantes en el sitio de la mol\'ecula central. Siempre apunta en la misma

direcci\'on que $\mathbf{E}$, lo que significa que el campo local es \emph{amplificado} por la
polarizaci\'on del medio. Esta amplificaci\'on es especialmente importante para materiales de
alta permitividad y es la raz\'on f\'isica por la que la relaci\'on de Clausius--Mossotti tiene la
forma que tiene.

\subsection{Relaci\'on de Clausius--Mossotti}

Sustituyendo el campo local en la ecuaci\'on de polarizaci\'on:

\begin{equation}
P = N\alpha\left(E + \frac{P}{3\varepsilon_0}\right)
\quad\implies\quad
P\!\left(1 - \frac{N\alpha}{3\varepsilon_0}\right) = N\alpha E
\end{equation}

Comparando con la definici\'on $P = \varepsilon_0(\varepsilon_r - 1)E$, e igualando
ambas expresiones para $P$, se cancela el campo $E$ y se obtiene, tras
reorganizaci\'on algebraica, la relaci\'on de Clausius--Mossotti:

\begin{equation}
\boxed{
\frac{\varepsilon_r - 1}{\varepsilon_r + 2} = \frac{N\alpha}{3\varepsilon_0}
}
\label{eq:CM}
\end{equation}

Esta relaci\'on establece un puente fundamental entre la respuesta microsc\'opica de
los \'atomos, descrita por la polarizabilidad $\alpha$, y la respuesta macrosc\'opica
del material, representada por la permitividad relativa $\varepsilon_r$. F\'isicamente,
implica que la respuesta diel\'ectrica de un medio no es simplemente la suma de
dipolos independientes, sino que incluye efectos colectivos derivados de la
interacci\'on entre ellos. Este es uno de los pilares conceptuales del presente trabajo.

\subsubsection{Significado F\'isico de la Relaci\'on de Clausius--Mossotti}

Esta relaci\'on establece un puente fundamental entre la respuesta microsc\'opica de
los \'atomos, descrita por la polarizabilidad $\alpha$, y la respuesta macrosc\'opica del
material, representada por la permitividad relativa $\varepsilon_r$. F\'isicamente,
implica que la respuesta diel\'ectrica de un medio no es simplemente la suma de
dipolos independientes, sino que incluye efectos colectivos derivados de la interacci\'on
entre ellos. Este es uno de los pilares conceptuales del presente trabajo.

La Ec.~\eqref{eq:CM} es el puente fundamental entre el mundo microsc\'opico ($\alpha$ de un
\'atomo individual) y el mundo macrosc\'opico ($\varepsilon_r$ del material en volumen).

\begin{itemize}
\item El lado izquierdo, $(\varepsilon_r - 1)/(\varepsilon_r + 2)$, es una funci\'on no
lineal de $\varepsilon_r$. El denominador $(\varepsilon_r + 2)$ codifica la
amplificaci\'on autoconsistente de muchos cuerpos.
\item El lado derecho, $N\alpha/3\varepsilon_0$, es la densidad de polarizabilidad.
\item En el l\'imite diluido ($N\alpha \ll 3\varepsilon_0$): $\varepsilon_r \approx
1 + N\alpha/\varepsilon_0$ --- una suma simple y aditiva.
\item En el l\'imite denso: el denominador $(\varepsilon_r + 2)$ produce una amplificaci\'on
no lineal de la constante diel\'ectrica.

\end{itemize}

\subsection{Medio Efectivo: Aproximaci\'on de Maxwell--Garnett}

Para medios heterog\'eneos (como una suspensi\'on de c\'elulas en plasma), la permitividad
efectiva no es un promedio simple. Para inclusiones esf\'ericas de permitividad $\varepsilon_i$ en
un medio hu\'esped $\varepsilon_m$ con fracci\'on volum\'etrica $f$:

\begin{equation}
\varepsilon_{\text{ef}} =
\varepsilon_m\left(
\frac{\varepsilon_i + 2\varepsilon_m + 2f(\varepsilon_i - \varepsilon_m)}
{\varepsilon_i + 2\varepsilon_m - f(\varepsilon_i - \varepsilon_m)}
\right)
\label{eq:MG}
\end{equation}

Esta forma funcional es estructuralmente id\'entica a la relaci\'on de Clausius--Mossotti,
y con raz\'on: ambas expresan la respuesta de polarizaci\'on autoconsistente de una colecci\'on
de objetos inmersos en un medio. Se aplicar\'a la Ec.~\eqref{eq:MG} a la sangre en la
secci\'on~\ref{sec:sangre}.

%=================================================================

\section{Mec\'anica Estad\'istica y Sistemas Dipolares}

\label{sec:stat}

%=================================================================

\noindent\textit{La secci\'on anterior asumi\'o que la polarizaci\'on surge del desplazamiento de
cargas, sin grados de libertad orientacionales. Esto funciona bien para los electrones. Pero las
mol\'eculas con momentos dipolares permanentes ---como el agua y los grupos cabeza de los
l\'ipidos de la membrana celular--- tambi\'en pueden rotar. Su alineaci\'on con el campo es
continuamente contrarrestada por la agitaci\'on t\'ermica. Esta secci\'on deriva la respuesta de
equilibrio y din\'amica de tal polarizaci\'on orientacional.}
\subsection{Estad\'istica de Boltzmann: Polarizaci\'on de Equilibrio}

Un dipolo permanente de momento $p$ en campo $E$ tiene energ\'ia potencial $U = pE\cos\theta$.

La distribuci\'on de Boltzmann da:

\begin{equation}
P(\theta) \propto e^{\frac{pE\cos\theta}{k_BT}}
\end{equation}

La proyecci\'on promedio a lo largo del eje del campo es la funci\'on de Langevin:

\begin{equation}
\langle p_z\rangle =
\frac{\displaystyle\int_0^\pi p\cos\theta\,
e^{\frac{pE\cos\theta}{k_BT}}\sin\theta\,d\theta}
{\displaystyle\int_0^\pi
e^{\frac{pE\cos\theta}{k_BT}}\sin\theta\,d\theta}
\end{equation}

En el l\'imite de campo d\'ebil ($pE \ll k_BT$, siempre v\'alido en bioimpedancia):

\begin{equation}
\langle p_z\rangle = \frac{p^2}{3k_BT}\,E
\end{equation}

La polarizabilidad orientacional es proporcional a $p^2/(k_BT)$: la agitaci\'on t\'ermica
combate la alineaci\'on, y el aumento de temperatura reduce la respuesta.

\subsection{Ecuaci\'on de Fokker--Planck: Din\'amica Rotacional}

La din\'amica de la distribuci\'on orientacional $f(\theta,t)$ est\'a gobernada por la ecuaci\'on de
Fokker--Planck:

\begin{equation}
\frac{\partial f}{\partial t} =
D_r\,\frac{1}{\sin\theta}\frac{\partial}{\partial\theta}
\left[\sin\theta\left(
\frac{\partial f}{\partial\theta}
+ \frac{1}{k_BT}f\,\frac{\partial U}{\partial\theta}
\right)\right]
\end{equation}

Tomando el primer momento angular $\langle\cos\theta\rangle$:

\begin{equation}
\frac{d}{dt}\langle\cos\theta\rangle =
-\frac{1}{\tau}\left(\langle\cos\theta\rangle - \frac{pE}{3k_BT}\right),
\qquad \tau = \frac{1}{2D_r}
\end{equation}

Multiplicando por $Np$, se obtiene la ecuaci\'on macrosc\'opica de Debye:

\begin{equation}
\tau\frac{dP}{dt} + P = P_{\text{eq}} = \varepsilon_0\chi_e\,E
\end{equation}

El tiempo de relajaci\'on $\tau = 1/(2D_r) = 4\pi\eta a^3/(k_BT)$ (Stokes--Einstein) depende de la
viscosidad $\eta$ y el tama\~{n}o molecular $a$: sistemas m\'as grandes o m\'as viscosos
se relajan m\'as lentamente.

%=================================================================

\section{Polarizaci\'on Din\'amica: Modelos de Debye y Cole--Cole}

\label{sec:debye}

%=================================================================

\noindent\textit{Habiendo derivado la din\'amica de relajaci\'on a partir de la ecuaci\'on de
Fokker--Planck, ahora se puede escribir la permitividad en el dominio de la frecuencia. Este es el
puente directo entre el movimiento molecular y el espectro diel\'ectrico que mide la
bioimpedancia.}

\subsection{Relajaci\'on de Debye en el Dominio de la Frecuencia}
Tomando la transformada de Fourier de la ecuaci\'on de Debye:

\begin{equation}
\boxed{
\varepsilon^*(\omega) =
\varepsilon_\infty +
\frac{\varepsilon_s - \varepsilon_\infty}{1 + j\omega\tau}
}
\label{eq:debye}
\end{equation}

Esta ecuaci\'on describe la din\'amica de relajaci\'on de la polarizaci\'on
del sistema. El tiempo caracter\'istico $\tau$ representa la escala temporal en
la que los dipolos reorientan su direcci\'on para alinearse con el campo el\'ectrico.
F\'isicamente, este modelo implica que la respuesta del sistema no es instant\'anea,
sino que presenta un retardo asociado a procesos microsc\'opicos como rotaci\'on
molecular o reconfiguraci\'on estructural. Aqu\'i nace el concepto de \emph{memoria
del sistema}: la polarizaci\'on actual depende no solo del campo presente, sino de
la historia del sistema.

\subsubsection{Significado F\'isico de Cada Par\'ametro}

Esta ecuaci\'on describe la din\'amica de relajaci\'on de la polarizaci\'on del sistema.
El tiempo caracter\'istico $\tau$ representa la escala temporal en la que los dipolos

reorientan su direcci\'on para alinearse con el campo el\'ectrico. F\'isicamente, este
modelo implica que la respuesta del sistema no es instant\'anea, sino que presenta un
retardo asociado a procesos microsc\'opicos como rotaci\'on molecular o reconfiguraci\'on
estructural. Aqu\'i nace el concepto de \emph{memoria del sistema}: la polarizaci\'on
actual depende no solo del campo presente, sino de la historia din\'amica del sistema.

\begin{itemize}
\item $\varepsilon_s$: \emph{permitividad est\'atica} ($\omega \to 0$) --- los dipolos
tienen tiempo suficiente para alinearse completamente.
\item $\varepsilon_\infty$: \emph{permitividad de alta frecuencia} ($\omega \to \infty$)
--- solo contribuye la polarizaci\'on electr\'onica.
\item $\tau$: \emph{tiempo de relajaci\'on} --- la transici\'on ocurre en
$f_r = 1/(2\pi\tau)$.
\item $j\omega\tau$: el denominador complejo hace que $\varepsilon^*$ sea compleja
--- la polarizaci\'on est\'a fuera de fase y se disipa energ\'ia por ciclo.
\end{itemize}

\subsubsection{Aspecto del Espectro de Debye}

\begin{itemize}
\item $\varepsilon'(\omega) = \text{Re}(\varepsilon^*)$: decrece como un escal\'on
sigmoidal de $\varepsilon_s$ a $\varepsilon_\infty$, centrado en $1/\tau$.
\item $\varepsilon''(\omega) = -\text{Im}(\varepsilon^*)$: tiene un pico en $\omega =
1/\tau$, que representa la m\'axima absorci\'on de energ\'ia.
\end{itemize}
\subsection{Generalizaci\'on de Cole--Cole}

Los sistemas biol\'ogicos reales se desvían del modelo de Debye ideal porque contienen
poblaciones de mol\'eculas con diferentes tama\~{n}os, entornos y estados de ligadura, lo que
resulta en una \emph{distribuci\'on} de tiempos de relajaci\'on en lugar de un \'unico $\tau$.

El modelo de Cole--Cole da cuenta de esto:

\begin{equation}
\boxed{
\varepsilon^*(\omega) =
\varepsilon_\infty +
\frac{\varepsilon_s - \varepsilon_\infty}{1 + (j\omega\tau)^{1-\alpha_{CC}}}
}
\label{eq:colecole}
\end{equation}

donde $\alpha_{CC} \in [0,1]$ ensancha el espectro de relajaci\'on. En $\alpha_{CC} = 0$, la
Ec.~\eqref{eq:colecole} se reduce exactamente al modelo de Debye.

El par'ametro $\alpha_{CC}$ introduce una distribuci'on de tiempos de relajaci'on,
lo cual es caracter'istico de sistemas complejos como los biol'ogicos. A diferencia
del modelo de Debye ideal, donde existe un 'unico tiempo caracter'istico, en tejidos
reales coexisten m'ultiples procesos din'amicos. Esto refleja la heterogeneidad
estructural del sistema, incluyendo variaciones en tama\~{n}o celular, composici'on
de membranas y propiedades del medio extracelular. El valor de $\alpha_{CC}$ puede
considerarse un biomarcador de la heterogeneidad biol'ogica del tejido estudiado. En sistemas
biol\'ogicos,
$\alpha_{CC} \neq 0$ surge de la variabilidad en tama\~{n}o celular, la heterogeneidad de la
membrana y las estructuras intracelulares complejas.

La contribuci\'on conductiva de los iones libres se agrega por separado:

\begin{equation}
\varepsilon^*(\omega) = \varepsilon_{\text{CC}}(\omega) - \frac{j\sigma}{\omega}
\end{equation}

El t\'ermino $\sigma/\omega$ diverge a bajas frecuencias, dominando la parte imaginaria de $Z$ y
dando origen al componente \'ohmico (resistivo) de la bioimpedancia.

%=================================================================

\section{Fundamentos Biol\'ogicos y Electroqu\'imicos de la Membrana Celular}

\label{sec:bio_foundations}
%=================================================================

\noindent\textbf{Transici\'on del modelo electrodin\'amico al biol\'ogico.}
Hasta este punto, el modelo ha sido desarrollado desde una perspectiva puramente
electrodin\'amica: polarizaci\'on at\'omica, permitividad macrosc\'opica, relajaci\'on
diel\'ectrica e impedancia efectiva de suspensiones celulares. Sin embargo, en sistemas
biol\'ogicos reales, la respuesta el\'ectrica est\'a fuertemente determinada por la estructura
y din\'amica de la membrana celular, cuyo comportamiento trasciende el de un simple
diel\'ectrico pasivo. A continuaci\'on se introducen los fundamentos fisiol\'ogicos y
bioqu\'imicos que complementan el modelo electrodin\'amico, conectando la f\'isica con la
biolog\'ia celular.

\noindent\textit{Las secciones anteriores establecieron la respuesta diel\'ectrica de medios
homog\'eneos desde primeros principios. Antes de modelar la c\'elula completa, es necesario
comprender en detalle la estructura y el comportamiento electroqu\'imico de la membrana
plasm\'atica, que es la interfaz cr\'itica que gobierna la dispersi\'on $\beta$ y conecta la f\'isica
con la biolog\'ia celular.}

\subsection{Estructura de la Membrana Celular}

La membrana celular es una estructura biol\'ogica compuesta principalmente por una bicapa
lip\'idica, formada por mol\'eculas de fosfolip\'idos anfif\'ilicos. Cada l\'ipido consiste en:

\begin{itemize}
\item Una cabeza hidrof\'ilica (polar)
\item Colas hidrof\'obicas (no polares)
\end{itemize}

Esta disposici\'on conduce a la formaci\'on de una bicapa, donde las colas hidrof\'obicas apuntan
hacia el interior y las cabezas polares interact\'uan con el entorno acuoso. Incrustadas en esta
estructura se encuentran prote\'inas, incluyendo canales i\'onicos, transportadores y receptores,
que regulan el intercambio de iones y mol\'eculas a trav\'es de la membrana.

Desde una perspectiva f\'isica, esta estructura genera una separaci\'on espacial entre regiones
de alta y baja permitividad:

\begin{itemize}
\item Regiones polares (cabezas): alta permitividad debido a interacciones dipolares
\item N\'ucleo hidrof\'obico (colas): baja permitividad debido a la ausencia de cargas libres

\end{itemize}

Como resultado, la membrana se comporta como una losa diel\'ectrica de espesor $d \sim
5\,\text{nm}$, separando dos medios conductores (citoplasma y fluido extracelular). Esta
discontinuidad diel\'ectrica es esencial para la separaci\'on de cargas y el confinamiento del
campo el\'ectrico a trav\'es de la membrana.

Las membranas celulares est\'an compuestas principalmente por bicapas lip\'idicas,
cuya estructura anf\'ip\'atica introduce una separaci\'on de cargas y una respuesta
diel\'ectrica altamente no homog\'enea. Esta caracter\'istica es fundamental para
la formaci\'on de potenciales de membrana y para la respuesta capacitiva que
domina el comportamiento de la bioimpedancia en la regi\'on de dispersi\'on $\beta$.
La membrana celular no solo act\'ua como un diel\'ectrico pasivo, sino como un
sistema din\'amico cuya respuesta depende de la apertura y cierre de canales
i\'onicos espec\'ificos, principalmente de sodio y potasio.

En este sentido, la membrana celular act\'ua como una barrera diel\'ectrica altamente selectiva.
Esta estructura introduce una capacitancia efectiva $C_m$ que regula el flujo de corriente i\'onica
entre el interior y el exterior celular, y cuya respuesta dependiente de la frecuencia constituye el
origen microsc\'opico de la dispersi\'on $\beta$ observada en los espectros de bioimpedancia.

\subsection{Propiedades El\'ectricas de la Membrana}

Desde una perspectiva electrodin\'amica, la membrana se comporta como una capa diel\'ectrica
con baja conductividad y permitividad finita. Debido a la separaci\'on de medios conductores
(fluidos intracelular y extracelular), la membrana act\'ua como un capacitor, almacenando
carga en sus interfaces. La capacitancia de membrana por unidad de \'area es:

\begin{equation}
C_m = \frac{\varepsilon_m}{d}
\end{equation}

donde $\varepsilon_m$ es la permitividad de la membrana y $d$ su espesor.

\subsubsection{Derivaci\'on de la Capacitancia desde la Ley de Gauss}

Considerando la membrana como un capacitor de placas paralelas de espesor $d$ y \'area $A$,
desde la ley de Gauss:

\begin{equation}
\oint \mathbf{D} \cdot d\mathbf{A} = Q
\end{equation}

Aplicando la ley de Gauss a una superficie cilíndrica que envuelve una de las
interfaces de la membrana, el vector de desplazamiento el\'ectrico $\mathbf{D}$
solo contribuye en la cara perpendicular a la membrana. Para un campo uniforme
dentro del diel\'ectrico de permitividad $\varepsilon_m$:

\begin{equation}
D = \varepsilon_m E = \frac{Q}{A}
\end{equation}

Despejando el campo el\'ectrico como $E = Q/(\varepsilon_m A)$ e integrando a lo
largo del espesor $d$ de la membrana, se obtiene la diferencia de potencial
entre las dos caras:

\begin{equation}
V_m = E\,d = \frac{Q\,d}{\varepsilon_m A}
\end{equation}

Reordenando esta expresi\'on para obtener la raz\'on carga--voltaje, se llega a la
capacitancia de membrana por unidad de \'area:

\begin{equation}
\boxed{C_m = \frac{Q}{V_m} = \frac{\varepsilon_m A}{d}}
\end{equation}
Esta expresi\'on conecta directamente con el modelo microsc\'opico: la permitividad
$\varepsilon_m$ proviene de la polarizabilidad at\'omica a trav\'es de la relaci\'on de
Clausius--Mossotti, estableciendo el enlace:

\begin{equation}
\alpha \rightarrow \varepsilon \rightarrow C_m
\end{equation}

\subsection{Interpretaci\'on Electrodin\'amica de la Membrana}

La membrana puede modelarse como una interfaz diel\'ectrica que separa dos regiones
conductoras con diferentes propiedades el\'ectricas. Bajo un campo el\'ectrico aplicado, las
condiciones de
frontera imponen:

\begin{equation}
\varepsilon_1^* E_1^\perp = \varepsilon_2^* E_2^\perp
\end{equation}

Esta discontinuidad conduce a la acumulaci\'on de cargas ligadas y libres en las interfaces,
generando polarizaci\'on interfacial (efecto Maxwell--Wagner). En este sentido, la membrana
existe en el modelo porque \emph{rompe} el campo el\'ectrico; su presencia es lo que genera
la dispersi\'on $\beta$.

\subsection{Potencial de Membrana}

El potencial de membrana surge del desequilibrio de concentraciones i\'onicas a trav\'es de la
membrana. Esto genera una diferencia de potencial el\'ectrico:

\begin{equation}
V_m = \phi_{\text{inside}} - \phi_{\text{outside}}
\end{equation}

Este potencial se mantiene por la permeabilidad selectiva de la membrana y los mecanismos de
transporte activo.

Desde la perspectiva electrodin\'amica desarrollada en este trabajo, la separaci\'on de cargas
est\'a directamente relacionada con los mecanismos de polarizaci\'on y los efectos capacitivos:

\begin{equation}
V_m = \frac{Q}{C_m}
\end{equation}

La carga $Q$ resulta tanto de cargas de polarizaci\'on (cargas ligadas) como de cargas i\'onicas
(cargas libres). Por tanto, el potencial de membrana no es puramente electroqu\'imico, sino que
tambi\'en contiene una contribuci\'on diel\'ectrica proveniente de los fen\'omenos de

polarizaci\'on: esta es una de las contribuciones conceptuales centrales del presente trabajo.

\subsection{Potencial de Nernst}

Para una \'unica especie i\'onica, el potencial de equilibrio est\'a dado por:

\begin{equation}
E_i =
\frac{RT}{z_i F}
\ln\left(
\frac{[i]_{\text{outside}}}{[i]_{\text{inside}}}
\right)
\end{equation}

donde:

\begin{itemize}
\item $R$ es la constante de los gases
\item $T$ es la temperatura
\item $z_i$ es la carga i\'onica
\item $F$ es la constante de Faraday
\end{itemize}

En equilibrio, el flujo neto de iones a trav\'es de la membrana es cero. Esto ocurre cuando la
fuerza el\'ectrica equilibra el gradiente qu\'imico:

\begin{equation}
z_i e E = \frac{d\mu}{dx}
\end{equation}

Resolviendo este balance se obtiene la ecuaci\'on de Nernst. Este resultado introduce energ\'ia
qu\'imica como complemento al modelo el\'ectrico.

\subsection{La Constante de Faraday y la Conductividad I\'onica}

La constante de Faraday, $F = 96{,}485\,\text{C/mol}$, es la carga el\'ectrica
total por mol de carga elemental. Su significado f\'isico es directo: si
$N_A = 6.022\times10^{23}$ iones monovalentes transportan cada uno una carga
$e = 1.602\times10^{-19}\,\text{C}$, entonces:
\begin{equation}
F = N_A \cdot e = 6.022\times10^{23} \times 1.602\times10^{-19} = 96{,}485\;\text{C/mol}
\end{equation}

Esta constante es la que convierte concentraciones i\'onicas (mol/m$^3$) en
densidades de carga el\'ectrica (C/m$^3$), conectando la qu\'imica con la
electrodin\'amica. En el contexto de la conductividad i\'onica, la corriente
producida por iones de especie $i$ con valencia $z_i$, movilidad $\mu_i$
(m$^2$/V$\cdot$s) y concentraci\'on $c_i$ (mol/m$^3$) es:
\begin{equation}
\sigma_i = z_i \cdot F \cdot \mu_i \cdot c_i
\end{equation}

Para un medio con m\'ultiples especies i\'onicas (sodio, potasio, cloruro), la
conductividad total es la suma de contribuciones:

\begin{equation}
\boxed{
\sigma = F \sum_i z_i \cdot \mu_i \cdot c_i
}
\label{eq:sigma_ionica}
\end{equation}

donde las movilidades el\'ectricas (a $25\,^\circ\text{C}$) son:
$\mu_{\text{Na}^+} = 5.2\times10^{-8}\,\text{m}^2/(\text{V}\cdot\text{s})$ y
$\mu_{\text{K}^+} = 7.6\times10^{-8}\,\text{m}^2/(\text{V}\cdot\text{s})$.

Para concentraciones fisiol\'ogicas t\'ipicas del citoplasma
($c_{\text{Na}^+} \approx 140\,\text{mM}$, $c_{\text{K}^+} \approx 5\,\text{mM}$,
con conversi\'on $1\,\text{mM} = 1\,\text{mol/m}^3$):
\begin{equation}
\sigma \approx F\left(\mu_{\text{Na}^+}\cdot 140 + \mu_{\text{K}^+}\cdot 5\right)
\approx 0.74\;\text{S/m}
\end{equation}

Este resultado es consistente con mediciones experimentales de la conductividad
intracitoplasm\'atica. N\'otese que $F$ aparece tambi\'en en la ecuaci\'on de
Nernst y en la de Goldman--Hodgkin--Katz, unificando la descripci\'on
electroqu\'imica de la membrana con el modelo de bioimpedancia: es la constante
que convierte procesos qu\'imicos (gradientes de concentraci\'on) en
fenomenolog\'ia el\'ectrica (potenciales y corrientes).

\subsection{Ecuaci\'on de Goldman--Hodgkin--Katz}

Cuando m\'ultiples iones est\'an involucrados, el potencial de membrana est\'a dado por:

\begin{equation}
V_m =
\frac{RT}{F}
\ln\left(
\frac{
P_K[K^+]_{\text{out}} + P_{Na}[Na^+]_{\text{out}} + P_{Cl}[Cl^-]_{\text{in}}
}{
P_K[K^+]_{\text{in}} + P_{Na}[Na^+]_{\text{in}} + P_{Cl}[Cl^-]_{\text{out}}
}
\right)
\end{equation}

donde $P_i$ representa la permeabilidad de la membrana para cada ion.

\subsection{Canales I\'onicos y Conductancia Din\'amica}

Los canales i\'onicos son estructuras prote\'icas que permiten el transporte selectivo
de iones a trav\'es de la membrana. Su apertura y cierre introducen una dependencia
temporal y no lineal en la conductividad del sistema, ya que el estado de los canales
depende del potencial de membrana. Esto implica que $\sigma$ debe interpretarse como
una funci\'on din\'amica del tiempo y del estado fisiol\'ogico de la c\'elula, y no
como una constante material. Las conductancias dependientes del tiempo se modelan como:

\begin{equation}
I_i = g_i(V_m, t)\,(V_m - E_i)
\end{equation}

donde:

\begin{itemize}
\item $g_i$ es la conductancia del ion $i$ (depende del estado del canal)
\item $E_i$ es el potencial de Nernst
\item $V_m - E_i$ es la fuerza impulsora
\end{itemize}

Esto introduce una componente no lineal y din\'amica en el comportamiento de la membrana.

\subsection{Papel del Entorno Acuoso}

Los sistemas biol\'ogicos son predominantemente acuosos, y el agua contribuye
significativamente a la respuesta diel\'ectrica. Su alta permitividad y relajaci\'on dipolar
influyen fuertemente en las mediciones de bioimpedancia (dispersi\'on $\gamma$, $f >
1\,\text{GHz}$). Este componente ya fue considerado en el modelo Cole--Cole, pero su papel en el
contexto
del entorno i\'onico es igualmente relevante para la conductividad efectiva del medio.

\subsection{Interacciones C\'elula-C\'elula}

Las c\'elulas no se comportan como sistemas aislados, sino que interact\'uan el\'ectricamente a
trav\'es del medio circundante. Estas interacciones modifican la impedancia efectiva e
introducen comportamiento colectivo, lo que es capturado por la teor\'ia de Maxwell--Garnett

en la secci\'on~\ref{sec:sangre}.

\subsection{Conexi\'on con el Modelo Electrodin\'amico}

El potencial de membrana puede interpretarse como una consecuencia macrosc\'opica de los
procesos de polarizaci\'on microsc\'opica. Al conectar las propiedades diel\'ectricas con el
comportamiento electroqu\'imico, se vuelve posible inferir estados fisiol\'ogicos a partir de
mediciones puramente el\'ectricas. La cadena conceptual de este trabajo se extiende entonces a:

\begin{equation}
\alpha \rightarrow \varepsilon^*(\omega) \rightarrow C_m \rightarrow V_m \rightarrow
I_{\text{ion}}
\end{equation}
%=================================================================

\section{Modelado de la Membrana Biol\'ogica}

\label{sec:membrane}

%=================================================================

\noindent\textit{Esta secci\'on modela las propiedades el\'ectricas de la membrana plasm\'atica
en el contexto del modelo electrodin\'amico completo.}

\subsection{Estructura y Propiedades El\'ectricas de la Membrana}

La membrana plasm\'atica es una bicapa lip\'idica de $\sim5\,\text{nm}$ de grosor. Sus
propiedades el\'ectricas est\'an dominadas por el n\'ucleo hidrof\'obico, que tiene:

\begin{itemize}
\item Permitividad muy baja: $\varepsilon_2 \approx 5\varepsilon_0$
(cadenas hidrocarbonadas, sin grupos polares).
\item Conductividad extremadamente baja: $\sigma_2 \sim 10^{-7}\,\text{S/m}$
(esencialmente impermeable a los iones).
\end{itemize}

La capacitancia por unidad de \'area es:

\begin{equation}
C_m = \frac{\varepsilon_2}{d}
\approx \frac{5 \times 8.854\times10^{-12}}{5\times10^{-9}}
\approx 8.9\times10^{-3}\,\text{F/m}^2 \approx 0.89\,\mu\text{F/cm}^2
\end{equation}

Este valor cl\'asico de $\sim 1\,\mu\text{F/cm}^2$ para membranas biol\'ogicas fue medido
por primera vez por Fricke (1925) \citep{Fricke1925} y sigue siendo un valor de referencia
en biof\'isica de membranas.

\subsection{La Membrana como Barrera de Corriente Dependiente de la Frecuencia}

La impedancia de la membrana por unidad de \'area es:

\begin{equation}
Z_m(\omega) = \frac{1}{j\omega C_m}
\end{equation}

\subsubsection{Interpretaci\'on F\'isica a Diferentes Frecuencias}

\begin{itemize}
\item \textbf{Baja frecuencia}: $Z_m \gg R_{\text{citoplasma}}$; la membrana es un
circuito abierto; \emph{ninguna corriente entra a la c\'elula}; toda la corriente
fluye por el espacio extracelular. La impedancia medida refleja \'unicamente el
plasma.
\item \textbf{Alta frecuencia}: $Z_m \to 0$; la membrana queda en cortocircuito;
\emph{la corriente fluye libremente hacia el citoplasma}; la impedancia medida
incluye ambos compartimentos.
\end{itemize}

Esta transici\'on es la base f\'isica de la dispersi\'on $\beta$ y es la raz\'on por la que la
bioimpedancia puede sondear las propiedades intracelulares de forma no invasiva.

\noindent\textbf{Interpretaci\'on f\'isica profunda de la transici\'on en frecuencia.}

A bajas frecuencias, la capacitancia de membrana domina y bloquea el flujo de corriente
hacia el interior celular, de modo que toda la corriente circula por el espacio extracelular.

A altas frecuencias, el campo el\'ectrico penetra el citoplasma al cortocircuitar la

capacitancia; en este r\'egimen, la conductividad intracelular es la que determina la
respuesta. Esta transici\'on es la huella electr\'onica de la bicapa lip\'idica en el espectro
de impedancia.

\subsection{Permitividad Compleja de Cada Regi\'on}

A cada compartimento se le asigna una permitividad compleja dependiente de la frecuencia:

\begin{equation}
\varepsilon_i^*(\omega) = \varepsilon_i(\omega) - \frac{j\sigma_i}{\omega}
\end{equation}
En este contexto, $\varepsilon_i^*(\omega)$ representa la permitividad efectiva de la regi\'on $i$,
la cual incluye contribuciones tanto diel\'ectricas como conductivas. La parte real $\varepsilon_i$
captura el almacenamiento diel\'ectrico.

%=================================================================

\section{Modelo Electrodin\'amico de una C\'elula Individual}

\label{sec:cell}

%=================================================================

\noindent\textit{Con las propiedades diel\'ectricas de cada regi\'on definidas, ahora se resuelve el
problema electrost\'atico para la c\'elula completa.}

\subsection{Modelo Geom\'etrico: Tres Regiones Conc\'entricas}

Se modela la c\'elula biol\'ogica como una estructura esf\'erica conc\'entrica:

\begin{itemize}
\item \textbf{Regi\'on I} --- Citoplasma ($r < a$): $\varepsilon_1^*$, $\sigma_1$
\item \textbf{Regi\'on II} --- Membrana ($a < r < a+d$): $\varepsilon_2^*$, $\sigma_2$
\item \textbf{Regi\'on III} --- Medio extracelular ($r > a+d$): $\varepsilon_3^*$, $\sigma_3$
\end{itemize}

La geometr\'ia esf\'erica es una idealizaci\'on de la forma de disco bic\'oncavo de los eritrocitos,
justificada porque preserva la f\'isica esencial (polarizaci\'on interfacial, apantallamiento
de la membrana) mientras sigue siendo anal\'iticamente tratable.

\subsection{Ecuaci\'on Gobernante: Ecuaci\'on de Laplace Generalizada}

Bajo la aproximaci\'on cuasiestática:

\begin{equation}
\nabla\cdot(\varepsilon_i^*\nabla\Phi_i) = 0 \quad\text{en cada regi\'on}
\end{equation}

La ecuaci'on de Laplace en coordenadas esf'ericas admite soluciones separables en
arm'onicos esf'ericos de orden $\ell$. Para un campo aplicado uniforme $E_0$ a lo
largo del eje $z$, solo el modo dipolar ($\ell = 1$) contribuye, pues los modos

superiores se anulan por simetr'ia. Las soluciones regulares en cada regi'on son:
\begin{equation}
\Phi_1 = -A_1 r\cos\theta, \quad
\Phi_2 = -\!\left(A_2 r + \frac{B_2}{r^2}\right)\!\cos\theta, \quad
\Phi_3 = -\!\left(E_0 r - \frac{B_3}{r^2}\right)\!\cos\theta
\end{equation}

\subsection{Condiciones de Frontera}

En cada interfaz se satisfacen simult\'aneamente dos condiciones:

\begin{enumerate}
\item \textbf{Continuidad del potencial}: $\Phi_i = \Phi_{i+1}$ (potencial un\'ivoco).
\item \textbf{Continuidad del desplazamiento normal}:
$\varepsilon_i^*\,\partial\Phi_i/\partial r =
\varepsilon_{i+1}^*\,\partial\Phi_{i+1}/\partial r$ (continuidad de carga).
\end{enumerate}

Estas cuatro condiciones determinan los coeficientes $A_1, A_2, B_2, B_3$.

\subsection{Aproximaci\'on de Membrana Delgada y Permitividad Interior Efectiva}

Para $d \ll a$ ($d/a \approx 10^{-3}$), la membrana act\'ua como una impedancia superficial
efectiva y el problema de tres regiones se reduce a:

\begin{equation}
\boxed{
\varepsilon_{\text{int}}^*(\omega) =
\frac{\varepsilon_1^*(\omega)}{\displaystyle
1 + \frac{d}{a}\cdot\frac{\varepsilon_1^*(\omega)}{\varepsilon_2^*(\omega)}}
}
\label{eq:epsint}
\end{equation}

\noindent Esta expresi\'on, derivada por Pauly y Schwan~(1959) para una esfera
con capa delgada, es dimensionalmente consistente: el factor de correcci\'on
$d\varepsilon_1^*/(a\varepsilon_2^*)$ es el producto de dos razones
adimensionales.
A bajas frecuencias $\varepsilon_{\mathrm{int}}^*\!\to\!0$ (membrana
bloquea la corriente); a altas frecuencias
$\varepsilon_{\mathrm{int}}^*\!\to\!\varepsilon_1^*$ (membrana transparente).
El radio celular $a$ entra expl\'icitamente, haciendo el promedio Monte Carlo
sobre la distribuci\'on de tama\~nos f\'isicamente significativo.

\subsubsection{Significado F\'isico de la Ec.~\eqref{eq:epsint}}

\begin{itemize}
\item \textbf{Baja frecuencia} ($Z_m \gg 1/\varepsilon_1^*$):
$\varepsilon_{\text{int}}^* \approx 0$. El citoplasma aparece el\'ectricamente
invisible: la membrana lo apantalla completamente.
\item \textbf{Alta frecuencia} ($Z_m \to 0$): $\varepsilon_{\text{int}}^* \to
\varepsilon_1^*$. La membrana es transparente; el citoplasma queda completamente
expuesto.
\end{itemize}

Este apantallamiento dependiente de la frecuencia es el mecanismo microsc\'opico de la
dispersi\'on $\beta$.

\subsection{Correcciones Num\'ericas y Consistencia F\'isica}
\label{sec:numerical_corrections}

El modelo num\'erico fue refinado para garantizar consistencia f\'isica y
estabilidad computacional. Se introdujeron las siguientes correcciones:

\begin{itemize}
\item \textbf{F\'ormula de membrana delgada (Pauly--Schwan):}
La permitividad interior efectiva incorpora el radio celular $a$,
haciendo la expresi\'on adimensionalmente consistente. El cociente
$d\varepsilon_1^*/(a\varepsilon_2^*)$ reproduce correctamente el apantallamiento
a baja frecuencia y la transparencia a alta frecuencia.

\item \textbf{Regularizaci\'on en frecuencia:}
El t\'ermino $\sigma/\omega$ diverge cuando $\omega\to 0$. Se evita mediante
$\omega\leftarrow\omega+\delta$, $\delta=10^{-20}\,\text{rad/s}\ll 1$.

\item \textbf{Distribuci\'on log-normal de radios:}
Garantiza valores siempre positivos y representa m\'as fielmente la
variabilidad biol\'ogica observada en eritrocitos.

\item \textbf{Maxwell--Garnett completo:}
Para $\phi=0.40$ se requiere la forma completa de la Ec.~\eqref{eq:MG},
que incorpora interacciones c\'elula-c\'elula a fracciones volum\'etricas elevadas.

\item \textbf{Promedio sobre impedancias:}
Se promedia $Z_i(\omega)$ de cada c\'elula en lugar de las permitividades,
preservando la naturaleza no lineal del problema.

\item \textbf{Representaci\'on compleja para ML:}
El vector de caracter\'isticas contiene $\mathrm{Re}(Z)$ e $\mathrm{Im}(Z)$,
preservando la informaci\'on de fase esencial para recuperar $\tau$ y $C_m$.

\item \textbf{Singularidades en Hodgkin--Huxley:}
Las funciones de compuerta $\alpha_m(V)$ y $\alpha_n(V)$ presentan formas
$0/0$ en $V=-40\,\text{mV}$ y $V=-55\,\text{mV}$ respectivamente, resueltas
mediante la regla de L'H\^opital.
\end{itemize}

\subsection{Relajaci\'on de Maxwell--Wagner y Polarizabilidad Celular}

La polarizabilidad efectiva de la c\'elula es:

\begin{equation}
\alpha_{\text{c\'elula}} =
4\pi\varepsilon_3^* a^3
\frac{\varepsilon_{\text{int}}^* - \varepsilon_3^*}
{\varepsilon_{\text{int}}^* + 2\varepsilon_3^*}
\end{equation}

La frecuencia de transici\'on est\'a determinada por el tiempo de relajaci\'on de Maxwell-Wagner. Este tiempo surge debido a la acumulaci'on de carga en las interfaces entre
medios con diferentes propiedades el'ectricas, como ocurre entre la membrana celular
y el medio extracelular. F'isicamente, describe la capacidad del sistema para
polarizarse a nivel de estructura celular, siendo responsable de la dispersi'on
tipo $\beta$ observada en tejidos biol'ogicos:

\begin{equation}
\tau_{MW} =
\frac{\varepsilon_{\text{int}} + 2\varepsilon_3}{\sigma_{\text{int}} + 2\sigma_3}
\label{eq:tauMW}
\end{equation}

Este tiempo de relajaci\'on determina la frecuencia de corte de la dispersi\'on
$\beta$: por debajo de $f_{MW} = 1/(2\pi\tau_{MW})$ la membrana bloquea la corriente;
por encima, queda en cortocircuito y la corriente penetra el citoplasma. F\'isicamente,
$\tau_{MW}$ encapsula simult\'aneamente la geometr\'ia celular, la capacitancia de
membrana y la conductividad del citoplasma, convirti\'endose en el par\'ametro
diagn\'ostico central de la dispersi\'on $\beta$ en tejidos biol\'ogicos.

\begin{table}[H]
\centering

\caption{Tres reg\'imenes de frecuencia de la respuesta electrodin\'amica celular.}
\label{tab:regimenes}
\begin{tabular}{@{}llll@{}}
\toprule
\textbf{R\'egimen} & \textbf{Frecuencia} & \textbf{Descripci\'on f\'isica}
& \textbf{Elemento dominante} \\
\midrule
Bajo & $f \ll f_{MW}$ & Membrana bloquea corriente; c\'elula aislante
& $Z_m \gg R_{\text{cit}}$ \\
Transici\'on & $f \approx f_{MW}$ & Redistribuci\'on de carga interfacial
& $\tau_{MW}$ \\
Alto & $f \gg f_{MW}$ & Membrana transparente; citoplasma conduce
& $\sigma_1$, $\varepsilon_1$ \\
\bottomrule
\end{tabular}
\end{table}

%=================================================================

\section{Sistemas Biol\'ogicos: Sangre y Suspensiones Celulares}

\label{sec:sangre}

%=================================================================

\noindent\textit{Una c\'elula individual no es un objeto cl\'inicamente medible. Lo que se mide en
la pr\'actica es la respuesta el\'ectrica colectiva de millones de c\'elulas suspendidas en plasma: la
sangre.}

\subsection{Permitividad Efectiva de Maxwell--Wagner en Sangre}

En este contexto, $\sigma$ representa la conductividad efectiva de cada
compartimento, la cual incluye contribuciones i\'onicas asociadas al movimiento
de iones de sodio (Na$^+$), potasio (K$^+$) y cloro (Cl$^-$) en el medio
acuoso. Para el plasma extracelular, $\sigma_3 \approx 1.2\,\text{S/m}$;
para el citoplasma, $\sigma_1 \approx 0.5\,\text{S/m}$.

Para una suspensi\'on de c\'elulas con fracci\'on volum\'etrica $\phi$ (hematocrito), aplicando la
teor\'ia de Maxwell--Garnett con permitividades complejas:

\begin{equation}
\varepsilon_{\text{ef}}^* =
\varepsilon_3^*\left(
\frac{\varepsilon_{\text{c\'elula}}^* + 2\varepsilon_3^* +
2\phi(\varepsilon_{\text{c\'elula}}^* - \varepsilon_3^*)}
{\varepsilon_{\text{c\'elula}}^* + 2\varepsilon_3^* - \phi(\varepsilon_{\text{c\'elula}}^* - \varepsilon_3^*)}
\right)
\label{eq:MWH}
\end{equation}

Aumentar $\phi$ reduce la conductividad efectiva a bajas frecuencias y acerca la permitividad de
alta frecuencia al valor del citoplasma. Un cambio del 10\% en $\phi$ (como en la anemia
moderada) produce un desplazamiento medible en $|Z(\omega)|$ a lo largo de todo el rango de
la dispersi\'on $\beta$.

%=================================================================

\section{Bioimpedancia Macrosc\'opica}

\label{sec:macro}

%=================================================================

\subsection{Expresi\'on General de la Impedancia}

\begin{equation}
\boxed{
Z(\omega) = \frac{1}{\sigma_{\text{ef}}(\omega) + j\omega\,\varepsilon_{\text{ef}}^*(\omega)}
}
\label{eq:Z}
\end{equation}

La impedancia representa la respuesta global del sistema ante una excitaci\'on
el'ectrica. En este modelo, no es 'unicamente una magnitud el'ectrica, sino
una funci'on que encapsula la din'amica interna del sistema biol'ogico.
Los dos componentes de $Z$ llevan informaci'on f'isica distinta:

\begin{itemize}
\item $\text{Re}(Z) = R(\omega)$: \textbf{resistencia} --- energ\'ia disipada por conducci\'on
\'ohmica i\'onica.
\item $\text{Im}(Z) = -X(\omega)$ (con $X > 0$): \textbf{reactancia} --- energ\'ia almacenada
por la capacitancia de la membrana y la polarizaci\'on diel\'ectrica.

\end{itemize}

\subsection{Hip\'otesis geom\'etrica}

Se asume una geometr\'ia efectiva homog\'enea en la cual el
campo el\'ectrico puede considerarse aproximadamente
uniforme, permitiendo una descripci\'on \emph{lumped}.
Esta hip\'otesis es v\'alida cuando:
\begin{equation}
\lambda = \frac{c}{f\sqrt{\varepsilon_r}} \gg a
\end{equation}
condici\'on que se cumple para $f < 1\,\text{GHz}$ y
$a \sim 5\,\mu\text{m}$.
Sin esta hip\'otesis, $Z$ como escalar requerir\'ia
un an\'alisis de valor de frontera completo.

\subsection{Ecuaci\'on Maestra del Sistema}

\label{sec:master_equation}

La formulaci\'on completa del modelo puede resumirse en una \emph{ecuaci\'on maestra} que
unifica todos los mecanismos f\'isicos en una sola expresi\'on:

\begin{equation}
\boxed{
Z(\omega,t) =
\frac{1}{
\sigma(V_m(t)) +
j\omega
\left[
\varepsilon_\infty +
\frac{\varepsilon_s - \varepsilon_\infty}
{1 + (j\omega\tau)^{1-\alpha_{CC}}}
\right]
}
}
\label{eq:master}
\end{equation}

con la conductividad modulada din\'amicamente por el potencial de membrana:

\begin{equation}
\sigma(V_m) = \sigma_0 \left(1 + \beta_\sigma\, V_m(t)\right)
\label{eq:sigma_vm}
\end{equation}
y $V_m(t)$ obtenido del sistema de Hodgkin--Huxley. Esta \'ultima ecuaci\'on
describe la din\'amica del potencial de membrana como un balance de corrientes:
la capacitiva (lado izquierdo), la externa aplicada $I_{\text{ext}}$, y las
i\'onicas de sodio, potasio y la fuga (lado derecho):

\begin{equation}
C_m \frac{dV_m}{dt} = I_{\text{ext}} - g_{Na}\,m^3 h\,(V_m - E_{Na})
- g_K\,n^4\,(V_m - E_K) - g_L\,(V_m - E_L)
\end{equation}

donde $m$, $h$ y $n$ son las variables de compuerta de los canales de Na$^+$ y K$^+$,
cuyos valores entre 0 y 1 modulan din\'amicamente la conductancia efectiva
de cada canal en funci\'on del voltaje y el tiempo.

\noindent La Ec.~\eqref{eq:master} es la s\'intesis de toda la cadena multiescala de este
trabajo: el numerador es la unidad, y el denominador encapsula \emph{simult\'aneamente} la
conducci\'on i\'onica modulada por el estado electrofisiol\'ogico de la c\'elula (t\'ermino

$\sigma(V_m)$) y la relajaci\'on diel\'ectrica distribuida de la membrana y el citoplasma (t\'ermino
de Cole--Cole). La dependencia en $t$ hace que $Z$ sea un observable din\'amico que evoluciona
con el potencial de acci\'on, mientras que la dependencia en $\omega$ lo hace
espectrosc\'opico. Ninguna formulaci\'on anterior ced\'{i}a ambas dependencias
simult\'aneamente
desde primeros principios.

\subsection{Separaci\'on de escalas: $\omega$ y $t$}
\label{sec:separacion_escalas}

El uso simult\'aneo de $\varepsilon^*(\omega)$ y
$\sigma(t)$ se justifica por la separaci\'on de escalas:
\begin{itemize}
\item Dependencia en $\omega$: respuesta lineal del
sistema alrededor del estado din\'amico instant\'aneo.
\item Dependencia en $t$: evoluci\'on lenta de
$\sigma(V_m(t))$ respecto a la excitaci\'on arm\'onica.
\end{itemize}
\begin{equation}
\tau_\text{fisiol\'ogico} \sim 1\,\text{ms}
\;\gg\;
\frac{1}{f_\text{medici\'on}} \sim 1\,\mu\text{s}
\end{equation}
Esta separaci\'on valida la aproximaci\'on cuasi-estacionaria
y el tratamiento conjunto de $\omega$ y $t$.

\paragraph{Interpretaci\'on del rol de cada par\'ametro.}

\begin{itemize}
\item $\tau$: controla la velocidad de relajaci\'on diel\'ectrica; desplaza la frecuencia de
transici\'on de la dispersi\'on $\beta$ seg\'un $f_{MW} \propto 1/(2\pi\tau)$.
\item $\alpha_{CC}$: refleja la heterogeneidad estructural de la membrana; valores mayores
corresponden a poblaciones celulares m\'as diversas.
\item $\varepsilon_s - \varepsilon_\infty$: amplitud de la dispersi\'on; proporcional a la
densidad de dipolos orientables.
\item $\sigma_0$, $\beta_\sigma$: conectan la excitabilidad neuronal (estado de canales
i\'onicos) con la conductividad macrosc\'opica medible.
\end{itemize}

\subsection{Cantidades Observables}
La observable primaria en este trabajo es el espectro de impedancia el\'ectrica:

\begin{equation}
Z(\omega)
\end{equation}

que codifica informaci\'on sobre la capacitancia de membrana, el transporte i\'onico y los
procesos de polarizaci\'on. Las observables secundarias incluyen el potencial de membrana
$V_m(t)$ y la permitividad efectiva $\varepsilon^*(\omega)$.

\subsection{Dualidad Frecuencia--Tiempo}

La respuesta de impedancia en el dominio de la frecuencia est\'a relacionada con la din\'amica
en el dominio del tiempo a trav\'es de transformadas de Fourier:

\begin{equation}
Z(\omega) \longleftrightarrow V_m(t)
\end{equation}

Esto establece un puente entre espectros de impedancia medibles y la din\'amica subyacente
de la membrana.

\subsection{Efecto de la Distribuci\'on de Tama\~{n}os Celulares}

Los radios celulares se muestrean de una distribuci\'on log-normal, garantizando siempre valores positivos y representando mejor la variabilidad biol\'ogica:

\begin{equation}
a \sim \mathrm{Log\text{-}Normal}(\ln\mu_a,\,0.2), \qquad \mu_a = 5\,\mu\text{m},\;
\sigma_a = 1\,\mu\text{m}
\end{equation}

La introducci\'on de una distribuci\'on de tama\~{n}os celulares permite capturar
la variabilidad biol\'ogica inherente a sistemas reales. Esta heterogeneidad es
responsable del ensanchamiento de la dispersi\'on $\beta$ observado en espectros
experimentales, el cual no puede reproducirse con un modelo unicelular de radio
fijo. Para cada radio muestreado se calcula $\varepsilon_{\text{c\'elula},i}^*(\omega)$
usando el modelo de membrana delgada. La permitividad efectiva de la suspensi\'on
se obtiene como el promedio de ensamble sobre las $N$ realizaciones:

\begin{equation}
\varepsilon_{\text{ef}}^*(\omega) =
\frac{1}{N}\sum_{i=1}^{N}\varepsilon_{\text{c\'elula},i}^*(\omega)
\end{equation}

Este promedio de conjunto ensancha la dispersi\'on $\beta$, explicando por qu\'e los espectros
experimentales siempre aparecen m\'as anchos que las predicciones unicelulares.
\noindent\textbf{El puente multiescala como fen\'omeno continuo.}
La relevancia de esta construcci\'on radica en que cada nivel de descripci\'on no es
independiente, sino que constituye una manifestaci\'on emergente del nivel anterior:
la polarizaci\'on at\'omica genera la permitividad, la permitividad determina la
impedancia de membrana, la impedancia modula el potencial y el potencial regula las
corrientes i\'onicas. As\'i, la bioimpedancia macrosc\'opica puede entenderse como una
consecuencia directa de procesos microsc\'opicos de polarizaci\'on. Este es el n\'ucleo
multiescala del presente trabajo.

\subsection{Formulaci\'on de la Impedancia Multiscala Completa}

Integrando todos los efectos descritos, la impedancia completa puede expresarse como:

\begin{equation}
Z(\omega,t) =
\left[
\int
\frac{1}{
\displaystyle
\sum_i z_i \mu_i c_i(t,\theta)
+
j\omega
\left(
\varepsilon_\infty(\theta) +
\frac{\varepsilon_s(\theta) - \varepsilon_\infty(\theta)}
{1 + (j\omega\tau(\theta))^{1-\alpha(\theta)}}
\right)
}
\,P(\theta)\,d\theta
\right]
\end{equation}

donde $\theta$ representa el conjunto de par\'ametros fisiol\'ogicos del sistema y $P(\theta)$
es la distribuci\'on de probabilidad de dichos par\'ametros.

Esta ecuaci\'on unifica din\'amica de cargas microsc\'opica, estructura celular mesosc\'opica
y respuesta el\'ectrica macrosc\'opica en una \'unica formulaci\'on multiscala de impedancia.

%=================================================================

\noindent\textbf{Conexi\'on final del marco te\'orico.}
En conjunto, este marco permite entender la bioimpedancia como una manifestaci\'on
macrosc\'opica de procesos f\'isicos microsc\'opicos, mediados por la estructura y
din\'amica celular. Esta perspectiva abre la posibilidad de utilizar mediciones
el\'ectricas como una herramienta para inferir propiedades fisiol\'ogicas internas.
La cadena multiescala completa es:
\begin{equation}
\underbrace{\alpha(\omega)}_{\text{\'atomo}}
\;\xrightarrow{\text{C-M}}\;
\underbrace{\varepsilon^*(\omega)}_{\text{medio}}
\;\xrightarrow{\text{M-W}}\;
\underbrace{Z_{\text{c\'el}}(\omega)}_{\text{c\'elula}}
\;\xrightarrow{\text{M-G}}\;
\underbrace{Z_{\text{sangre}}(\omega)}_{\text{tejido}}
\;\xrightarrow{\text{ML}}\;
\underbrace{\hat{\theta}}_{\text{par\'ametros}}
\end{equation}

%=================================================================

\section{Implementaci\'on Num\'erica}
\label{sec:numerical}

%=================================================================

\noindent\textit{Habiendo derivado el marco anal\'itico completo, esta secci\'on presenta su
realizaci\'on num\'erica. La simulaci\'on calcula el espectro completo de impedancia a partir de
par\'ametros f\'isicos y genera los datos de entrenamiento para el aprendizaje autom\'atico.}

\subsection{Malla de Frecuencias y Pipeline}

La simulaci\'on cubre $f \in [10^2, 10^9]\,\text{Hz}$ con $N_f = 600$ puntos equiespaciados
logar\'itmicamente, capturando los reg\'imenes de dispersi\'on $\alpha$, $\beta$ y $\gamma$
con igual resoluci\'on por d\'ecada.

\begin{enumerate}
\item Calcular $\varepsilon_i^*(\omega)$ para cada regi\'on (Cole--Cole + conductividad).
\item Calcular $Z_m(\omega)$ y $\varepsilon_{\text{int}}^*(\omega)$ (modelo de membrana
delgada).
\item Extraer $N = 100$ radios celulares de $\mathcal{N}(\mu_a, \sigma_a)$; promediar
$\varepsilon_{\text{c\'elula}}^*(\omega)$.
\item Calcular $Z(\omega)$.
\end{enumerate}

\subsection{Organizaci\'on del C\'odigo en M\'odulos}

La implementaci\'on se estructura en cuatro archivos Python independientes y ejecutables
en secuencia, siguiendo la cadena f\'isica del modelo:

\begin{center}
\renewcommand{\arraystretch}{1.35}
\begin{tabular}{@{}lll@{}}
\toprule
\textbf{Archivo} & \textbf{Contenido} & \textbf{Salida} \\
\midrule
\texttt{01\_simulacion.py}     & Pipeline f\'isico completo; barrido de par\'ametros & PDFs de espectros \\
\texttt{02\_generar\_dataset.py} & Generador Monte Carlo; base de datos SQLite & \texttt{bioimpedancia.db} \\
\texttt{03\_entrenar\_red.py}  & Entrenamiento BioNet; normalizaci\'on; curva de p\'erdida & \texttt{bionet\_inverso.pt} \\
\texttt{04\_resultados.py}     & Gr\'aficas de validaci\'on; m\'etricas MAE, MAPE, $R^2$ & PDFs de resultados \\
\bottomrule
\end{tabular}
\end{center}

El orden de ejecuci\'on es estrictamente secuencial:
\begin{enumerate}
\item \texttt{python 01\_simulacion.py} \hfill\textit{$\sim$10 segundos}
\item \texttt{python 02\_generar\_dataset.py} \hfill\textit{$\sim$5--10 minutos}
\item \texttt{python 03\_entrenar\_red.py} \hfill\textit{$\sim$2--5 minutos}
\item \texttt{python 04\_resultados.py} \hfill\textit{$\sim$30 segundos}
\end{enumerate}

\subsection{C\'odigo de Simulaci\'on Completo}

\begin{lstlisting}[language=Python, caption={Simulaci\'on completa
(corregida): Pauly--Schwan + Maxwell--Garnett completo + promedio de poblaci\'on.}]
import numpy as np
import matplotlib.pyplot as plt

f     = np.logspace(2, 9, 600)   # 100 Hz a 1 GHz
omega = 2 * np.pi * f
omega = omega + 1e-20            # [C1] evitar division por cero
epsilon0 = 8.854e-12

def cole_cole(eps_inf, eps_s, tau, alpha, omega):
    return eps_inf + (eps_s - eps_inf) / (1 + (1j*omega*tau)**(1 - alpha))

eps1_inf = 5*epsilon0;  eps1_s = 60*epsilon0
tau1 = 1e-6;  alpha1 = 0.1;  sigma1 = 0.5       # Citoplasma
eps2 = 5*epsilon0;  sigma2 = 1e-7;  d = 5e-9    # Membrana
eps3_inf = 10*epsilon0; eps3_s = 80*epsilon0
tau3 = 1e-8;  alpha3 = 0.05;  sigma3 = 1.2      # Plasma
phi = 0.4                                         # Hematocrito

eps1 = cole_cole(eps1_inf, eps1_s, tau1, alpha1, omega) - 1j*sigma1/omega
eps3 = cole_cole(eps3_inf, eps3_s, tau3, alpha3, omega) - 1j*sigma3/omega
eps2_complex = eps2 - 1j*sigma2/omega

def simular_poblacion(n_celulas=100):
    # [C5] Log-normal: radios siempre positivos
    radios = np.random.lognormal(mean=np.log(5e-6), sigma=0.2, size=n_celulas)
    Z_total = np.zeros_like(omega, dtype=complex)
    for a in radios:
        # [C2] Pauly-Schwan 1959 (usa radio a; dimensionalmente correcto)
        eps_int = eps1 / (1 + (d/a)*(eps1/eps2_complex))
        # [C3] Maxwell-Garnett completo (valido para phi=0.4)
        eps_cell = eps3 * (
            (eps_int + 2*eps3 + 2*phi*(eps_int - eps3)) /
            (eps_int + 2*eps3 - phi*(eps_int - eps3))
        )
        # [C4] Promediar IMPEDANCIAS (preserva no linealidad)
        Z_total += 1 / (sigma3 + 1j*omega*eps_cell)
    return Z_total / n_celulas

# [C6] CORRECCION BUG CRITICO:
# simular_poblacion() ya devuelve la impedancia promediada.
# Incorrecto (original): eps_ef = simular_poblacion(100)
#                        Z = 1/(sigma3+jw*eps_ef)  <- formula aplicada 2 veces
# Correcto:
Z = simular_poblacion(100)

fig, ejes = plt.subplots(1, 2, figsize=(13, 5))
ejes[0].loglog(f, np.abs(Z))
ejes[0].set(xlabel='Frecuencia (Hz)', ylabel=r'$|Z|$ ($\Omega$)',
            title='Magnitud de Impedancia')
ejes[0].grid(True, which='both', alpha=0.4)
ejes[1].plot(np.real(Z), -np.imag(Z))
ejes[1].set(xlabel=r'Re$(Z)$', ylabel=r'$-$Im$(Z)$',
            title='Diagrama de Nyquist')
ejes[1].grid(True, alpha=0.4)
plt.tight_layout()
plt.savefig('bioimpedancia_espectros.pdf', dpi=150, bbox_inches='tight')
plt.show()
\end{lstlisting}

\subsection{Validaci\'on de l\'imites f\'isicos del c\'odigo}

La implementaci\'on num\'erica reproduce correctamente los
comportamientos asint\'oticos conocidos:
\begin{itemize}
\item $\omega \to 0$: $Z \approx 1/\sigma$
(conducci\'on i\'onica dominante).
\item $\omega \to \infty$: $|Z| \propto 1/\omega$
(respuesta diel\'ectrica).
\end{itemize}
Esta verificaci\'on es obligatoria para garantizar que
los datos sint\'eticos del dataset sean f\'isicamente
coherentes antes de entrenar el modelo de ML.

\begin{lstlisting}[language=Python, caption={Funci\'on de validaci\'on de l\'imites f\'isicos (ejecutar antes de generar el dataset).}]
def validar_fisica(Z, omega, sigma_ref):
    # Verifica limites analiticos obligatorios.
    # Limite omega -> 0: Z ~ 1/sigma (conduccion ionica)
    # Limite omega -> inf: |Z| ~ 1/(w*eps) (dielectrico)
    Z_low  = np.abs(Z[0])
    Z_high = np.abs(Z[-1])
    print("Validacion de limites fisicos:")
    print(f"  |Z| baja f: {Z_low:.3e}  | esperado ~{1/sigma_ref:.3e}")
    print(f"  |Z| alta f: {Z_high:.3e}")
    assert Z_high < Z_low, "ERROR: Z no decrece con frecuencia"
    assert np.all(np.isfinite(Z)), "ERROR: Z contiene NaN o Inf"
    print("  OK: todos los limites fisicos verificados.")
\end{lstlisting}

\subsection{Implementaci\'on del Modelo F\'isico con Conductividad I\'onica}

El siguiente c\'odigo integra la conductividad i\'onica derivada de par\'ametros fisiol\'ogicos:

\begin{lstlisting}[language=Python, caption={Modelo f\'isico completo con conductividad i\'onica y capacitancia.}]
import numpy as np

f     = np.logspace(2, 9, 500)
omega = 2*np.pi*f + 1e-20       # [C1] regularizacion
epsilon0 = 8.854e-12

def cole_cole(eps_inf, eps_s, tau, alpha, omega):
    return eps_inf + (eps_s - eps_inf)/(1 + (1j*omega*tau)**(1-alpha))

def ionic_conductivity(c_na, c_k):
    # sigma = F * sum_i z_i*mu_i*c_i  (Ec. sigma_ionica)
    F           = 96485.0
    z_na, z_k   = 1, 1
    mu_na, mu_k = 5.2e-8, 7.6e-8
    return F * (z_na*mu_na*c_na + z_k*mu_k*c_k)

eps_inf = 5  * epsilon0
eps_s   = 60 * epsilon0
tau     = 1e-6
alpha   = 0.1
c_na    = 140.0;  c_k = 5.0
sigma   = ionic_conductivity(c_na, c_k)
eps_complex = cole_cole(eps_inf, eps_s, tau, alpha, omega) - 1j*sigma/omega

A  = 1e-10;  d = 5e-9
Cm = eps_complex * A / d
g_Na = 1e-6;  g_K = 2e-6;  G = g_Na + g_K
Z       = 1 / (G + 1j*omega*Cm)
Z_mag   = np.abs(Z)
Z_phase = np.angle(Z, deg=True)
\end{lstlisting}

\subsection{Implementaci\'on Monte Carlo: Promedio Sobre Distribuci\'on de Par\'ametros}

La implementaci\'on num\'erica sigue un enfoque Monte Carlo para aproximar la impedancia
promediada sobre el ensemble, capturando la variabilidad biol\'ogica mediante muestreo
estoc\'astico de par\'ametros fisiol\'ogicos:

\begin{lstlisting}[language=Python, caption={Implementaci\'on Monte Carlo de la impedancia multiscala.}]
import numpy as np

f     = np.logspace(2, 9, 500)
omega = 2*np.pi*f + 1e-20
epsilon0 = 8.854e-12

def cole_cole(eps_inf, eps_s, tau, alpha, omega):
    return eps_inf + (eps_s - eps_inf)/(1 + (1j*omega*tau)**(1-alpha))

def ionic_conductivity(c_na, c_k):
    F = 96485.0;  z_na, z_k = 1, 1
    mu_na, mu_k = 5.2e-8, 7.6e-8
    return F * (z_na*mu_na*c_na + z_k*mu_k*c_k)

def single_realization(omega):
    c_na  = np.random.normal(140, 10)
    c_k   = np.random.normal(5, 1)
    tau   = np.random.uniform(1e-7, 1e-4)
    alpha = np.random.uniform(0.05, 0.3)
    eps_inf = 5 * epsilon0;  eps_s = 80 * epsilon0
    sigma   = ionic_conductivity(c_na, c_k)
    eps_c   = cole_cole(eps_inf, eps_s, tau, alpha, omega) - 1j*sigma/omega
    return 1 / (sigma + 1j * omega * eps_c)

def monte_carlo_Z(omega, N=200):
    Z_total = np.zeros_like(omega, dtype=complex)
    for _ in range(N):
        Z_total += single_realization(omega)
    return Z_total / N

Z = monte_carlo_Z(omega, N=300)
\end{lstlisting}

%=================================================================

\section{Base de Datos e Integraci\'on de Aprendizaje Autom\'atico}
\label{sec:ml}

%=================================================================

\noindent\textit{La simulaci\'on f\'isica puede ejecutarse para cualquier combinaci\'on de
par\'ametros dentro de los rangos fisiol\'ogicos. Esta secci\'on usa esa capacidad para construir un
conjunto de datos etiquetados y entrenar redes neuronales para resolver el problema inverso. El
aprendizaje autom\'atico no se emplea como una caja negra, sino como el inverso matem\'atico
del modelo f\'isico.}

\subsection{Esquema de Base de Datos y Generaci\'on de Datos}

Cada registro de la base de datos almacena el espectro de impedancia (100 puntos de
frecuencia) como \emph{vector de caracter\'isticas} y los par\'ametros f\'isicos
$(\sigma_1, \sigma_3, \phi)$ como \emph{etiquetas}. El esquema SQL es:

\begin{lstlisting}[language=SQL, caption={Esquema de base de datos para bioimpedancia.}]
CREATE TABLE bioimpedance (
    id      INTEGER PRIMARY KEY,
    g_Na    REAL,
    g_K     REAL,
    tau     REAL,
    alpha   REAL,
    sigma   REAL,
    Cm      REAL,
    Z_real  BLOB,
    Z_imag  BLOB
);
\end{lstlisting}

\begin{lstlisting}[language=Python, caption={Generador Monte Carlo y base de datos SQLite (con correcciones f\'isicas).}]
import numpy as np, sqlite3, pickle
epsilon0 = 8.854e-12
f = np.logspace(3, 8, 100)

def simular_celula(f, params):
    e1, s1, e2, s2, e3, s3, d, a, phi = params
    omega  = 2 * np.pi * f + 1e-20          # [C1]
    eps1   = e1 - 1j*s1/omega
    eps3   = e3 - 1j*s3/omega
    eps2_c = e2 - 1j*s2/omega               # membrana compleja
    # [C2] Pauly-Schwan con radio a (dimensionalmente correcto)
    eps_int = eps1 / (1 + (d/a)*(eps1/eps2_c))
    # [C3] Maxwell-Garnett completo
    eps_ef = eps3 * (
        (eps_int + 2*eps3 + 2*phi*(eps_int - eps3)) /
        (eps_int + 2*eps3 - phi*(eps_int - eps3))
    )
    Z = 1 / (s3 + 1j * omega * eps_ef)
    # [C4] Re+Im (no solo magnitud)
    X = np.concatenate([np.real(Z), np.imag(Z)])
    return (X - np.mean(X)) / (np.std(X) + 1e-30)

conn = sqlite3.connect('bioimpedancia.db')
c    = conn.cursor()
c.execute("""CREATE TABLE IF NOT EXISTS simulaciones
    (id INTEGER PRIMARY KEY, sigma1 REAL, sigma3 REAL,
     phi REAL, d REAL, a REAL, espectro_z BLOB)""")
conn.commit()

N_MUESTRAS = 5000   # aumentar a 10000 para mayor precision
for _ in range(N_MUESTRAS):
    s1  = np.random.uniform(0.2, 1.0)
    s3  = np.random.uniform(0.5, 2.0)
    ph  = np.random.uniform(0.1, 0.6)
    d_  = np.random.uniform(3e-9, 8e-9)
    a_  = np.random.uniform(3e-6, 8e-6)
    params = (60*epsilon0, s1, 5*epsilon0, 1e-7, 80*epsilon0, s3, d_, a_, ph)
    X = simular_celula(f, params)
    c.execute('INSERT INTO simulaciones (sigma1,sigma3,phi,d,a,espectro_z) '
              'VALUES (?,?,?,?,?,?)', (s1, s3, ph, d_, a_, pickle.dumps(X)))
    if (_ + 1) % 500 == 0:
        conn.commit()   # commit cada 500 muestras
conn.commit();  conn.close()
\end{lstlisting}

\subsection{Modelos Directo e Inverso}

\begin{equation}
\text{Directo:}\; (\tau,\sigma,\varepsilon)
\xrightarrow{\text{RN}_D} \hat{Z}(\omega)
\qquad
\text{Inverso:}\; Z(\omega)
\xrightarrow{\text{RN}_I} (\hat{\tau},\hat{\sigma},\hat{\varepsilon})
\end{equation}

El modelo directo (RN$_D$) sirve como validaci\'on: reproduce el espectro a partir
de par\'ametros conocidos. El modelo inverso (RN$_I$) es el de mayor relevancia
cl\'inica: infiere par\'ametros fisiol\'ogicos a partir de la medici\'on.
Adicionalmente, una tercera red clasifica el estado fisiol\'ogico global:

\begin{equation}
\text{Clasificaci\'on:}\; Z(\omega)
\xrightarrow{\text{RN}_C} \text{estado fisiol\'ogico}
\end{equation}

\subsection{Por qu\'e el Aprendizaje Autom\'atico Funciona}

\label{sec:why_ml}

La red neuronal tiene \'exito en el problema inverso por una raz\'on f\'isica profunda: el mapeo
entre espectros de impedancia y par\'ametros fisiol\'ogicos es \emph{determin\'istico dentro del
modelo f\'isico}. No hay ambig\"uedad estoc\'astica en la generaci\'on de datos; cada vector de
caracter\'isticas corresponde a exactamente un conjunto de par\'ametros. Esto convierte el
problema inverso en una regresi\'on bien definida que la red puede aprender con precisi\'on.

Si los datos de entrenamiento provinieran de fuentes emp\'iricas o ad-hoc, la red aprender\'ia
correlaciones espurias sin base f\'isica y no podr\'ia generalizar. La f\'isica es, en este
sentido, el garante de la generalizaci\'on del modelo de aprendizaje autom\'atico.

\begin{center}
\textit{
La red neuronal tiene \'exito precisamente porque el mapeo entre espectros de impedancia
y par\'ametros fisiol\'ogicos es f\'isicamente restringido y determin\'istico dentro del
modelo propuesto --- no es simplemente un ajuste de curvas.
}
\end{center}

\noindent\textbf{Naturaleza del problema inverso.}
El problema inverso consiste en determinar los par'ametros f'isicos del sistema a
partir de su respuesta en frecuencia. Este problema es intr'insecamente mal
condicionado, ya que m'ultiples configuraciones pueden producir respuestas similares,
particulamente cuando $\sigma$ y $\tau$ var'ian simult'aneamente. El uso de redes
neuronales se justifica precisamente por la naturaleza no lineal y altamente acoplada
de este problema, el cual resulta dif'icil de resolver anal'iticamente.

\noindent\textbf{Nota sobre la interpretabilidad del modelo de aprendizaje autom\'atico.}
Es importante destacar que el modelo de aprendizaje autom\'atico no sustituye la f\'isica
subyacente, sino que act\'ua como una herramienta para aproximar la soluci\'on inversa del
sistema. Su capacidad de generalizaci\'on est\'a garantizada por el hecho de que los datos
de entrenamiento provienen \'integramente del modelo f\'isico, lo cual asegura que la red
aprende relaciones f\'isicamente fundamentadas y no correlaciones emp\'iricas superficiales.

\subsection{Arquitectura y Entrenamiento de la Red Neuronal (Modelo Inverso)}
La red mapea un espectro de impedancia de 200 caracter\'isticas
$\mathbf{x} \in \mathbb{R}^{200}$ (100 puntos de $\text{Re}(Z)$ + 100 de $\text{Im}(Z)$)
a los cinco par\'ametros fisiol\'ogicos objetivo
$\mathbf{y} = (g_{Na},\,g_K,\,\sigma_{\text{base}},\,\tau,\,\alpha_{CC})$.

\textbf{Elecci\'on de la arquitectura.} Se seleccion\'o una red totalmente conectada
(200\,$\to$\,256\,$\to$\,128\,$\to$\,64\,$\to$\,5 neuronas)
con \emph{Batch Normalization} antes de cada ReLU.
El cuello de botella 256\,$\to$\,128\,$\to$\,64 fuerza a la red
a aprender una representaci\'on comprimida del espectro.
La incorporaci\'on de BatchNorm estabiliza el entrenamiento cuando los
par\'ametros de entrada abarcan varios \'ordenes de magnitud~\citep{Barsoukov2005}.

\textbf{Funci\'on de p\'erdida.} MSE con normalizaci\'on $z$-score de salidas
(cr\'itico: $g_{Na}\sim 100$ vs.\ $\tau\sim 10^{-5}$, escalas muy diferentes).

\textbf{Regularizaci\'on.} Dropout ($p = 0.1$) y decaimiento de pesos ($\lambda = 10^{-4}$).
Programador de tasa de aprendizaje StepLR (factor $0.5$ cada 15 \'epocas).

\begin{lstlisting}[language=Python, caption={Red neuronal BioNet con BatchNorm para el problema inverso de bioimpedancia (5 par\'ametros fisiol\'ogicos).}]
import torch, torch.nn as nn, numpy as np, sqlite3, pickle
from torch.utils.data import Dataset, DataLoader, random_split

class BioDataset(Dataset):
    def __init__(self, db):
        conn = sqlite3.connect(db)
        filas = conn.execute(
            'SELECT g_Na, g_K, sigma_base, tau, alpha_cc, espectro_z '
            'FROM simulaciones').fetchall()
        conn.close()
        self.X = np.array([pickle.loads(f[5]) for f in filas], dtype=np.float32)
        self.y = np.array([[f[0],f[1],f[2],f[3],f[4]] for f in filas],
                          dtype=np.float32)
        # Normalizacion z-score de etiquetas [critico: escalas muy diferentes]
        self.y_mean = self.y.mean(axis=0)
        self.y_std  = self.y.std(axis=0) + 1e-30
        self.y = (self.y - self.y_mean) / self.y_std
    def __len__(self):  return len(self.X)
    def __getitem__(self, i):
        return (torch.tensor(self.X[i], dtype=torch.float32),
                torch.tensor(self.y[i], dtype=torch.float32))

class BioNet(nn.Module):
    """
    200 -> 256 -> 128 -> 64 -> 5
    BatchNorm antes de ReLU estabiliza el entrenamiento [C6].
    """
    def __init__(self, input_size=200, n_params=5):
        super().__init__()
        self.model = nn.Sequential(
            nn.Linear(input_size, 256),
            nn.BatchNorm1d(256), nn.ReLU(), nn.Dropout(0.1),
            nn.Linear(256, 128),
            nn.BatchNorm1d(128), nn.ReLU(), nn.Dropout(0.1),
            nn.Linear(128, 64),  nn.ReLU(),
            nn.Linear(64, n_params))
    def forward(self, x):  return self.model(x)

dataset = BioDataset('bioimpedancia.db')
n_val   = int(len(dataset) * 0.20)
train_ds, val_ds = random_split(dataset, [len(dataset)-n_val, n_val])
train_loader = DataLoader(train_ds, batch_size=64, shuffle=True)
val_loader   = DataLoader(val_ds,   batch_size=64)

modelo   = BioNet(input_size=dataset.X.shape[1])
optimiz  = torch.optim.Adam(modelo.parameters(), lr=1e-3, weight_decay=1e-4)
planif   = torch.optim.lr_scheduler.StepLR(optimiz, step_size=15, gamma=0.5)
perdida  = nn.MSELoss()

for epoca in range(50):
    modelo.train(); total = 0
    for X, y in train_loader:
        pred = modelo(X);  p = perdida(pred, y)
        optimiz.zero_grad();  p.backward();  optimiz.step()
        total += p.item()
    modelo.eval()
    with torch.no_grad():
        val_loss = perdida(modelo(torch.tensor(
            dataset.X[val_ds.indices], dtype=torch.float32)),
            torch.tensor(dataset.y[val_ds.indices], dtype=torch.float32)).item()
    planif.step()
    if (epoca+1) % 10 == 0:
        print(f'Epoca {epoca+1:3d} | Train: {total/len(train_loader):.6f} '
              f'| Val: {val_loss:.6f}')

torch.save({'model_state_dict': modelo.state_dict(),
            'y_mean': dataset.y_mean, 'y_std': dataset.y_std},
           'bionet_inverso.pt')
\end{lstlisting}

\subsection{Principio de Consistencia F\'isica del Dataset}

Una limitaci\'on cr\'itica en enfoques anteriores es la generaci\'on de datos de entrenamiento
con modelos simplificados o sint\'eticos que no reflejan la f\'isica real del sistema. Si el
dataset no proviene del modelo f\'isico completo, la red neuronal aprende ruido artificial
en lugar de relaciones f\'isicas reales.

En este trabajo, \emph{cada muestra del dataset se genera a partir del pipeline f\'isico
completo}:

\begin{equation}
\text{Par\'ametros fisiol\'ogicos}
\;\xrightarrow{\text{HH}}\;
V_m(t)
\;\xrightarrow{\sigma(V)}\;
\sigma(t)
\;\xrightarrow{\text{Cole--Cole}}\;
\varepsilon^*(\omega)
\;\xrightarrow{}\;
Z(\omega,t)
\;\xrightarrow{\text{features}}\;
\mathbf{x}
\end{equation}

\subsection{Par\'ametros Aleatorizados del Dataset}

\begin{center}
\renewcommand{\arraystretch}{1.3}
\begin{tabular}{@{}lll@{}}
\toprule
\textbf{Par\'ametro} & \textbf{Rango} & \textbf{Significado} \\
\midrule
$g_{Na}$ & $[80,\;140]$ & Conductancia m\'axima de Na$^+$ \\
$g_K$ & $[20,\;50]$ & Conductancia m\'axima de K$^+$ \\
$\sigma_{\text{base}}$ & $[0.5,\;2.0]\;\text{S/m}$ & Conductividad base \\
$\tau$ & $[10^{-6},\;10^{-4}]\;\text{s}$ & Tiempo de relajaci\'on \\
$\alpha$ & $[0.05,\;0.3]$ & Ensanchamiento Cole--Cole \\
\bottomrule
\end{tabular}
\end{center}

\subsection{Pipeline Completo: HH + Conductividad Din\'amica + Cole--Cole}

\begin{lstlisting}[language=Python, caption={Pipeline f\'isico completo: Hodgkin--Huxley + conductividad din\'amica + Cole--Cole.}]
import numpy as np
epsilon0 = 8.854e-12

def simulate_HH(g_Na, g_K, T=50, dt=0.02):
    t = np.arange(0, T, dt)
    V = -65*np.ones_like(t);  m = np.zeros_like(t)
    h = np.zeros_like(t);     n = np.zeros_like(t)
    E_Na, E_K, E_L = 50, -77, -54.4
    g_L = 0.3;  C_m = 1.0;  I_ext = 10.0
    def alpha_m(V):       # [C7] singularidad V=-40
        dV = V+40.0
        if abs(dV)<1e-7: return 1.0
        return 0.1*dV/(1.0-np.exp(-dV/10.0))
    def beta_m(V):  return 4*np.exp(-(V+65)/18)
    def alpha_h(V): return 0.07*np.exp(-(V+65)/20)
    def beta_h(V):  return 1/(1+np.exp(-(V+35)/10))
    def alpha_n(V):       # [C7] singularidad V=-55
        dV = V+55.0
        if abs(dV)<1e-7: return 0.1
        return 0.01*dV/(1.0-np.exp(-dV/10.0))
    def beta_n(V):  return 0.125*np.exp(-(V+65)/80)
    for i in range(len(t)-1):
        am,bm = alpha_m(V[i]),beta_m(V[i])
        ah,bh = alpha_h(V[i]),beta_h(V[i])
        an,bn = alpha_n(V[i]),beta_n(V[i])
        m[i+1] = m[i] + dt*(am*(1-m[i]) - bm*m[i])
        h[i+1] = h[i] + dt*(ah*(1-h[i]) - bh*h[i])
        n[i+1] = n[i] + dt*(an*(1-n[i]) - bn*n[i])
        I_Na = g_Na*(m[i]**3)*h[i]*(V[i]-E_Na)
        I_K  = g_K*(n[i]**4)*(V[i]-E_K)
        I_L  = g_L*(V[i]-E_L)
        V[i+1] = V[i] + dt*(I_ext - I_Na - I_K - I_L)/C_m
    return V

def sigma_dynamic(V, sigma_base):
    return sigma_base * (1 + 0.01*(V + 65))

def cole_cole(omega, eps_inf, eps_s, tau, alpha):
    return eps_inf + (eps_s-eps_inf)/(1+(1j*omega*tau)**(1-alpha))

def compute_Z(V, sigma_base, tau, alpha):
    f     = np.logspace(3, 7, 100)
    omega = 2*np.pi*f + 1e-20    # [C1]
    eps   = cole_cole(omega, 5*epsilon0, 80*epsilon0, tau, alpha)
    Z_time = []
    for v in V:
        sigma = sigma_dynamic(v, sigma_base)
        # eps = pure Cole-Cole (no conductivity term); sigma handled separately
        Z_time.append(1 / (sigma + 1j*omega*eps))
    return np.array(Z_time)

def extract_features(Z):
    # [C4] Re+Im solamente (no |Z|)
    Z_mean = np.mean(Z, axis=0)
    X = np.concatenate([np.real(Z_mean), np.imag(Z_mean)])
    return (X - X.mean()) / (X.std() + 1e-30)

def generate_physical_sample(noise_level=0.02):
    g_Na       = np.random.uniform(80, 140)
    g_K        = np.random.uniform(20, 50)
    sigma_base = np.random.uniform(0.5, 2.0)
    tau        = np.random.uniform(1e-6, 1e-4)
    alpha      = np.random.uniform(0.05, 0.3)
    V = simulate_HH(g_Na, g_K)
    Z = compute_Z(V, sigma_base, tau, alpha)
    Z = Z + noise_level * np.random.randn(*Z.shape)   # ruido 2%
    X = extract_features(Z)
    y = np.array([g_Na, g_K, sigma_base, tau, alpha])
    return X, y

X_data, y_data = [], []
for _ in range(1000):
    X, y = generate_physical_sample()
    X_data.append(X);  y_data.append(y)
X_data = np.array(X_data);  y_data = np.array(y_data)
\end{lstlisting}

\noindent\textbf{Nota sobre consistencia f\'isica.} Todos los datos de entrenamiento fueron
generados usando un modelo basado en f\'isica que combina la din\'amica de Hodgkin--Huxley
con la respuesta diel\'ectrica dependiente de la frecuencia. Al incluir parte real e imaginaria
(no solo magnitud), el vector de caracter\'isticas preserva la informaci\'on de fase que codifica
la capacitancia de membrana y los tiempos de relajaci\'on.

\subsection{Normalizaci\'on y Divisi\'on Train/Test}

La normalizaci\'on es cr\'itica cuando los par\'ametros objetivo abarcan diferentes escalas
f\'isicas ($g_{Na} \sim 100$ versus $\tau \sim 10^{-5}$):

\begin{lstlisting}[language=Python, caption={Normalizaci\'on z-score y divisi\'on 80/20 train/test.}]
from sklearn.model_selection import train_test_split

X_mean = X_data.mean(axis=0);  X_std = X_data.std(axis=0)
y_mean = y_data.mean(axis=0);  y_std = y_data.std(axis=0)
X_norm = (X_data - X_mean) / (X_std + 1e-30)
y_norm = (y_data - y_mean) / (y_std + 1e-30)

X_train, X_test, y_train, y_test = train_test_split(
    X_norm, y_norm, test_size=0.2, random_state=42)
\end{lstlisting}

\subsection{Arquitectura Mejorada: BioNet con BatchNorm}

\begin{lstlisting}[language=Python, caption={Arquitectura BioNet mejorada con BatchNorm para 5 par\'ametros.}]
import torch
import torch.nn as nn
from torch.utils.data import DataLoader, TensorDataset

class BioNet(nn.Module):
    def __init__(self, input_size, n_params=5):  # [C9] 5 parametros fisicos
        super().__init__()
        self.model = nn.Sequential(
            nn.Linear(input_size, 256),
            nn.BatchNorm1d(256),   # [C6] BatchNorm ANTES de ReLU
            nn.ReLU(),
            nn.Dropout(0.1),
            nn.Linear(256, 128),
            nn.BatchNorm1d(128),
            nn.ReLU(),
            nn.Dropout(0.1),
            nn.Linear(128, 64),
            nn.ReLU(),
            nn.Linear(64, n_params)  # [C9] 5 params fisicos: sigma1,sigma3,phi,d,a
        )
    def forward(self, x):  return self.model(x)

X_train_t = torch.tensor(X_train, dtype=torch.float32)
y_train_t = torch.tensor(y_train, dtype=torch.float32)
X_test_t  = torch.tensor(X_test,  dtype=torch.float32)
y_test_t  = torch.tensor(y_test,  dtype=torch.float32)

train_loader = DataLoader(TensorDataset(X_train_t, y_train_t),
                          batch_size=32, shuffle=True)
model     = BioNet(input_size=X_train.shape[1])
optimizer = torch.optim.Adam(model.parameters(), lr=1e-3)
loss_fn   = nn.MSELoss()

for epoch in range(30):
    model.train();  total_loss = 0
    for X_batch, y_batch in train_loader:
        pred = model(X_batch)
        loss = loss_fn(pred, y_batch)
        optimizer.zero_grad();  loss.backward();  optimizer.step()
        total_loss += loss.item()
    model.eval()
    with torch.no_grad():
        val_loss = loss_fn(model(X_test_t), y_test_t).item()
    print(f'Epoch {epoch+1:3d} | Train: {total_loss:.4f} | Val: {val_loss:.4f}')
\end{lstlisting}

\subsection{Ruido Experimental: Realismo del Dataset}

Un modelo entrenado con datos perfectos (sin ruido) puede fallar en datos experimentales
reales, donde siempre existe ruido de medici\'on instrumental. Para hacer el dataset
fisiol\'ogicamente realista y mejorar la robustez del modelo, se a\~{n}ade ruido gaussiano
calibrado al espectro de impedancia simulado:

\begin{lstlisting}[language=Python, caption={Adici\'on de ruido gaussiano al 2\% para realismo experimental.}]
def generate_physical_sample(noise_level=0.02):
    g_Na       = np.random.uniform(80, 140)
    g_K        = np.random.uniform(20, 50)
    sigma_base = np.random.uniform(0.5, 2.0)
    tau        = np.random.uniform(1e-6, 1e-4)
    alpha      = np.random.uniform(0.05, 0.3)
    V     = simulate_HH(g_Na, g_K)
    Z     = compute_Z(V, sigma_base, tau, alpha)
    # Ruido gaussiano ~2% SNR (Sanchez et al. 2013)
    noise = noise_level * np.random.randn(*Z.shape)
    Z_noisy = Z + noise
    X = extract_features(Z_noisy)
    y = np.array([g_Na, g_K, sigma_base, tau, alpha])
    return X, y
\end{lstlisting}

El nivel de ruido del 2\% es consistente con las incertidumbres t\'ipicas reportadas en
sistemas de espectroscop\'ia de bioimpedancia de cuatro electrodos \citep{Sanchez2013}.

Entrenar con datos ruidosos fuerza a la red a aprender caracter\'isticas robustas del
espectro en lugar de ajustar artefactos espec\'ificos de la simulaci\'on.

\subsection{An\'alisis de Sensibilidad: Correlaci\'on Par\'ametro--Espectro}

Antes de entrenar la red, es fundamental cuantificar qu\'e tan fuertemente est\'a codificado
cada par\'ametro en el espectro de impedancia. Un par\'ametro con baja correlaci\'on con las
caracter\'isticas ser\'a m\'as dif\'icil de recuperar por el modelo inverso. Este an\'alisis
tambi\'en identifica qu\'e partes del espectro son m\'as informativas para cada par\'ametro:

\begin{lstlisting}[language=Python, caption={An\'alisis de correlaci\'on entre caracter\'isticas del espectro y par\'ametros.}]
import numpy as np
import matplotlib.pyplot as plt

corr_matrix = np.corrcoef(X_data.T, y_data.T)
n_feat      = X_data.shape[1]
corr_Xy     = corr_matrix[:n_feat, n_feat:]
labels_y    = ['g_Na', 'g_K', 'sigma_base', 'tau', 'alpha']
max_corr    = np.max(np.abs(corr_Xy), axis=0)

print('Correlacion maxima |r|:')
for name, c in zip(labels_y, max_corr):
    print(f'  {name:12s}: {c:.3f}')

fig, ax = plt.subplots(figsize=(10, 4))
im = ax.imshow(np.abs(corr_Xy).T, aspect='auto', cmap='hot', vmin=0, vmax=1)
ax.set_yticks(range(5));  ax.set_yticklabels(labels_y)
ax.set_xlabel('Indice de feature (Re | Im)')
ax.set_title('Correlacion |r| entre features del espectro y parametros')
plt.colorbar(im, ax=ax, label='|r|')
plt.tight_layout()
plt.savefig('correlacion_features.pdf', dpi=150, bbox_inches='tight')
plt.show()
\end{lstlisting}

\begin{figure}[H]
\centering
\framebox[0.95\textwidth]{\rule{0pt}{3.5cm}
\centering Mapa de correlaci\'on $|r|$ entre features del espectro y par\'ametros --- marcador}
\caption{Mapa de calor de correlaci\'on $|r|$ entre features del espectro y par\'ametros
fisiol\'ogicos. Filas: par\'ametros $(g_{Na}, g_K, \sigma_{\text{base}}, \tau, \alpha_{CC})$;
columnas: features del espectro (\'indices 0--99 corresponden a $\text{Re}(Z)$,
\'indices 100--199 a $\text{Im}(Z)$). Regiones m\'as brillantes indican mayor
contenido de informaci\'on sobre cada par\'ametro. Generado por \texttt{04\_resultados.py}.}
\label{fig:corr}
\end{figure}

\subsection{Evaluaci\'on: MAE y Error Relativo por Par\'ametro}

\begin{lstlisting}[language=Python, caption={M\'etricas cuantitativas del modelo inverso con desnormalizaci\'on.}]
import numpy as np

model.eval()
with torch.no_grad():
    preds_norm = model(X_test_t).numpy()

preds_real = preds_norm * y_std + y_mean
y_real     = y_test    * y_std + y_mean
labels     = ['g_Na', 'g_K', 'sigma_base', 'tau', 'alpha']

mae = np.mean(np.abs(preds_real - y_real), axis=0)
print('MAE por parametro:')
for name, err in zip(labels, mae):
    print(f'  {name:12s}: {err:.4f}')

relative_error = np.mean(
    np.abs((preds_real - y_real) / (y_real + 1e-10)), axis=0)
print('Error relativo (MAPE):')
for name, err in zip(labels, relative_error):
    print(f'  {name:12s}: {err*100:.2f}%')
\end{lstlisting}

\subsection{Gr\'aficas de Validaci\'on: Predicho vs.\ Real}

\begin{lstlisting}[language=Python, caption={Script completo de validaci\'on \texttt{04\_resultados.py}: m\'etricas MAE, MAPE, $R^2$ y gr\'aficas de dispersi\'on.}]
import numpy as np, torch, torch.nn as nn, sqlite3, pickle, matplotlib.pyplot as plt

LABELS = ['g_Na', 'g_K', 'sigma_base', 'tau', 'alpha_cc']
COLORES = ['#2c7bb6','#d7191c','#1a9641','#f29000','#7b2d8b']

# -- Cargar modelo guardado ---------------------------------------------------
ckpt  = torch.load('bionet_inverso.pt', map_location='cpu')
model = BioNet(input_size=200, n_params=5)
model.load_state_dict(ckpt['model_state_dict'])
model.eval()
y_mean = ckpt['y_mean'];  y_std = ckpt['y_std']

# -- Cargar 20% de test desde SQLite -----------------------------------------
np.random.seed(42)
conn  = sqlite3.connect('bioimpedancia.db')
filas = conn.execute('SELECT g_Na,g_K,sigma_base,tau,alpha_cc,espectro_z '
                     'FROM simulaciones').fetchall()
conn.close()
idx   = np.random.permutation(len(filas))[:int(len(filas)*0.2)]
X_test = np.array([pickle.loads(filas[i][5]) for i in idx], dtype=np.float32)
y_test = np.array([[filas[i][j] for j in range(5)] for i in idx], dtype=np.float32)

# -- Prediccion y desnormalizacion -------------------------------------------
with torch.no_grad():
    y_pred = model(torch.tensor(X_test)).numpy() * y_std + y_mean

# -- Metricas ----------------------------------------------------------------
mae  = np.mean(np.abs(y_pred - y_test), axis=0)
mape = np.mean(np.abs((y_pred-y_test)/(np.abs(y_test)+1e-10)),axis=0)*100
r2   = np.array([1 - np.sum((y_test[:,i]-y_pred[:,i])**2) /
                    np.sum((y_test[:,i]-y_test[:,i].mean())**2)
                 for i in range(5)])
print(f"{'Param':>10}  {'MAE':>10}  {'MAPE(%)':>8}  {'R2':>6}")
for n,m,p,r in zip(LABELS,mae,mape,r2):
    print(f"{n:>10}  {m:>10.4f}  {p:>8.2f}%  {r:>6.4f}")

# -- Grafica predicho vs real ------------------------------------------------
fig, axes = plt.subplots(1, 5, figsize=(22, 4))
fig.suptitle('Modelo Inverso -- Predicho vs. Real', fontsize=13)
for i, ax in enumerate(axes):
    ax.scatter(y_test[:,i], y_pred[:,i], alpha=0.35, s=8,
               color=COLORES[i], rasterized=True)
    lims = [min(y_test[:,i].min(), y_pred[:,i].min()),
            max(y_test[:,i].max(), y_pred[:,i].max())]
    ax.plot(lims, lims, 'k--', linewidth=1.2)
    ax.set_title(f'{LABELS[i]}\n$R^2$={r2[i]:.3f}, MAPE={mape[i]:.1f}\%',
                 fontsize=9)
    ax.set_xlabel('Real', fontsize=9);  ax.set_ylabel('Predicho', fontsize=9)
    ax.grid(True, alpha=0.3)
plt.tight_layout()
plt.savefig('04a_predicho_vs_real.pdf', dpi=150, bbox_inches='tight')
plt.show()
\end{lstlisting}

\begin{figure}[H]
\centering
\framebox[0.95\textwidth]{\rule{0pt}{3.5cm}
\centering Gr\'aficas predicho vs.\ real para $(g_{Na},\,g_K,\,\sigma,\,\tau,\,\alpha_{CC})$ --marcador}
\caption{Validaci\'on del modelo inverso. Cada panel muestra la dispersi\'on entre
el valor real del par\'ametro y el estimado por BioNet sobre el 20\,\% del conjunto de prueba.
Los puntos cercanos a la l\'inea $y=x$ (negra discontinua) indican alta precisi\'on.
Se reportan $R^2$ y MAPE por par\'ametro. Los resultados se obtienen ejecutando
\texttt{04\_resultados.py}.}
\label{fig:scatter}
\end{figure}

\subsection{Identificabilidad de Par\'ametros}

\label{sec:identifiability}

No todos los par\'ametros fisiol\'ogicos son igualmente identificables a partir de datos de
impedancia, ya que algunas combinaciones de par\'ametros producen respuestas espectrales
similares. Este fen\'omeno, conocido como \emph{degeneraci\'on del espacio de par\'ametros},
es una limitaci\'on intr\'inseca del problema inverso.

En particular, un incremento simult\'aneo de $\sigma$ y $\tau$ puede generar espectros
casi indistinguibles en regiones de alta conductividad, limitando la recuperaci\'on
individual de ambos par\'ametros. Para mejorar la identificabilidad se recomienda incluir
caracter\'isticas espectrales expl\'icitas, como la frecuencia del m\'aximo de la parte
imaginaria $\text{Im}(Z)$, que es m\'as sensible a $\tau$ que la magnitud $|Z|$ sola.

\subsection{Interpretaci\'on F\'isica de los Errores por Par\'ametro}

No todos los par\'ametros fisiol\'ogicos son igualmente f\'aciles de recuperar a partir del
espectro de impedancia. La raz\'on no es num\'erica sino f\'isica: cada par\'ametro deja una
``huella'' diferente en el espectro, y algunas huellas son m\'as \'unicas y localizadas que
otras. A continuaci\'on se analiza cada par\'ametro:

\paragraph{$g_{Na}$ (conductancia de sodio).} La conductancia m\'axima de Na$^+$ se predice
con \emph{error relativamente bajo}. Esto se
debe a que $g_{Na}$ controla la amplitud y la velocidad de la despolarizaci\'on en el
modelo de Hodgkin--Huxley, lo que produce un cambio de gran amplitud y breve duraci\'on en
$\sigma(t)$. Esta excursi\'on r\'apida y pronunciada deja una firma distintiva en la parte real del
espectro promediado, especialmente en la banda $10^5$--$10^7\,\text{Hz}$ donde la
corriente penetra el citoplasma.

\paragraph{$g_K$ (conductancia de potasio).} $g_K$ se predice con error \emph{moderado}. A
diferencia de Na$^+$, los canales de K$^+$
se activan m\'as lentamente y su efecto principal es la duraci\'on de la repolarizaci\'on, no
su amplitud. En el espectro promediado temporalmente, este efecto se diluye, reduciendo
la distinguibilidad. Una posible mejora es incluir caracter\'isticas del dominio temporal
(por ejemplo, la anchura del pulso en $Z(t)$) en el vector de features.

\paragraph{$\sigma_{\text{base}}$ (conductividad base).}

La conductividad base se recupera con \emph{alta precisi\'on} porque afecta

\emph{uniformemente} todo el espectro de impedancia: un aumento de $\sigma$ desplaza toda
la curva $|Z(\omega)|$ hacia abajo. Esta influencia global hace que sea f\'acilmente
identificable incluso con ruido moderado.

\paragraph{$\tau$ (tiempo de relajaci\'on Cole--Cole).} El par\'ametro $\tau$ exhibe el
\emph{mayor error de predicci\'on}. La raz\'on f\'isica es su influencia no lineal sobre la respuesta

en frecuencia: $\tau$ desplaza la frecuencia de transici\'on de la dispersi\'on $\beta$, pero este
desplazamiento se vuelve menos
distinguible en reg\'imenes de alta conductividad donde la magnitud de la dispersi\'on se
comprime. Formalmente:

\begin{equation}
f_{MW} = \frac{1}{2\pi\tau} \propto \frac{\sigma_{\text{int}} + 2\sigma_3}
{\varepsilon_{\text{int}} + 2\varepsilon_3}
\end{equation}
Un aumento simult\'aneo de $\sigma$ y $\tau$ puede producir espectros similares, creando
degeneraci\'on en el espacio de par\'ametros. Para mejorar su recuperaci\'on se recomienda
incluir la \emph{frecuencia del m\'aximo de la parte imaginaria} como feature expl\'icita.

\paragraph{$\alpha_{CC}$ (ensanchamiento Cole--Cole).} $\alpha_{CC}$ se predice con
\emph{precisi\'on intermedia}. Controla el ancho del
semic\'irculo en el diagrama de Nyquist, pero su efecto es sutil cuando el ruido es
comparable a la depresi\'on angular. Se beneficia de incluir features de la parte imaginaria
del espectro, donde su influencia es m\'as pronunciada.

\subsection{Conclusi\'on: $Z(\omega,t)$ Contiene Informaci\'on Fisiol\'ogica Invertible}

Los resultados de este trabajo establecen un resultado te\'oricamente no trivial y
pr\'acticamente relevante:

\begin{equation}
\boxed{
Z(\omega,t) \;\text{ contiene informaci\'on fisiol\'ogica invertible}
}
\end{equation}

Esto significa que el mapeo completo, desde par\'ametros fisiol\'ogicos hasta
el espectro de impedancia y de regreso, es aproximadamente invertible. La cadena
es la siguiente: el modelo f\'isico genera el espectro a partir de los par\'ametros;
la red neuronal invierte esa relaci\'on y recupera los par\'ametros originales:

\begin{equation}
(g_{Na},\;g_K,\;\sigma,\;\tau,\;\alpha)
\;\xrightarrow{\text{modelo f\'isico}}\;
Z(\omega,t)
\;\xrightarrow{\text{red neuronal}}\;
(\hat{g}_{Na},\;\hat{g}_K,\;\hat{\sigma},\;\hat{\tau},\;\hat{\alpha})
\end{equation}

es \emph{aproximadamente invertible}: la red neuronal recupera con precisi\'on aceptable los
par\'ametros fisiol\'ogicos a partir del espectro de impedancia. Este resultado no es trivial
porque el modelo directo es altamente no lineal (combina ecuaciones diferenciales de HH con
permitividad compleja de Cole--Cole), y la inversi\'on de sistemas no lineales generalmente
no tiene soluci\'on anal\'itica cerrada.

La demostraci\'on de invertibilidad tiene tres implicaciones inmediatas:

\begin{enumerate}
\item \textbf{No se entren\'o solo una red.} Se demostr\'o que $Z(\omega,t)$ contiene

informaci\'on fisiol\'ogica suficiente para reconstruir las propiedades celulares que la
generaron.
\item \textbf{La bioimpedancia es un observable diagn\'ostico completo.} Una medici\'on
el\'ectrica no invasiva puede, en principio, sustituir mediciones invasivas de conductancias
i\'onicas o composici\'on sangu\'inea.
\item \textbf{El modelo f\'isico es esencial.} Sin el modelo f\'isico como generador de
datos de entrenamiento, la red neuronal carecer\'ia de una base de verdad f\'isica y no
podr\'ia generalizar a nuevos estados fisiol\'ogicos. La f\'isica y el aprendizaje autom\'atico
son inseparables en este enfoque.
\end{enumerate}

%=================================================================

\section{Extensi\'on Fisiol\'ogica y Bioqu\'imica}

\label{sec:physiological}

%=================================================================

\noindent\textit{Las secciones anteriores establecieron un marco electrodin\'amico completo,
desde la polarizaci\'on at\'omica hasta la bioimpedancia macrosc\'opica. Sin embargo, los sistemas
biol\'ogicos no son medios puramente pasivos. La respuesta el\'ectrica de las c\'elulas est\'a
fuertemente influenciada por procesos electroqu\'imicos, transporte i\'onico y din\'amica de
membrana. Esta secci\'on extiende el modelo f\'isico hacia una descripci\'on fisiol\'ogica.}

\subsection{Motivaci\'on para la Extensi\'on Fisiol\'ogica}

El marco electrodin\'amico desarrollado en las secciones anteriores proporciona una descripci\'on
f\'isica completa de la polarizaci\'on, la respuesta diel\'ectrica y la bioimpedancia. Sin embargo,
los sistemas biol\'ogicos exhiben comportamientos que no pueden ser capturados completamente
por modelos el\'ectricos pasivos. En particular, las c\'elulas vivas mantienen activamente
gradientes el\'ectricos a trav\'es de mecanismos de transporte i\'onico. Por tanto, se hace
necesario extender el modelo f\'isico hacia una descripci\'on fisiol\'ogica y bioqu\'imica, con el
fin de conectar las propiedades el\'ectricas medibles con la funci\'on biol\'ogica subyacente.

\subsection{La Membrana Celular como Interfaz Electrodin\'amica Activa}

En el modelo f\'isico, la membrana fue tratada como una capa diel\'ectrica pasiva. Sin embargo,
en los sistemas biol\'ogicos, la membrana es una interfaz activa que regula el transporte
i\'onico y mantiene gradientes el\'ectricos. Por tanto, la membrana debe entenderse no solo
como un capacitor, sino como un sistema din\'amico que acopla procesos el\'ectricos, qu\'imicos
y biol\'ogicos.

\subsection{Mecanismos de Transporte Activo}

Las c\'elulas mantienen activamente gradientes i\'onicos a trav\'es de bombas dependientes de
ATP, como la Na$^+$/K$^+$-ATPasa. Esto introduce una componente de no-equilibrio en el
sistema,
sosteniendo el potencial de membrana a lo largo del tiempo. Sin este mecanismo, el sistema
tender\'ia hacia el equilibrio pasivo y no representar\'ia c\'elulas vivas reales.

\subsection{Origen F\'isico del Potencial de Membrana}

El potencial de membrana emerge de la separaci\'on de cargas a trav\'es de la membrana. Desde
la perspectiva electrodin\'amica desarrollada en este trabajo, esta separaci\'on de cargas est\'a
directamente relacionada con los mecanismos de polarizaci\'on y los efectos capacitivos:

\begin{equation}
Q = C_m V_m
\end{equation}

Por tanto, el potencial de membrana puede interpretarse como una consecuencia macrosc\'opica
de los procesos de polarizaci\'on microsc\'opica.

\subsection{Gradientes Electroqu\'imicos y Equilibrio I\'onico}

Las c\'elulas mantienen gradientes de concentraci\'on de iones como Na$^+$, K$^+$ y Cl$^-$ a
trav\'es de la membrana. Estos gradientes generan potenciales de equilibrio descritos por la
ecuaci\'on de Nernst:

\begin{equation}
E_i = \frac{RT}{z_i F}
\ln\left(
\frac{[i]_{\text{out}}}{[i]_{\text{in}}}
\right)
\end{equation}

Esto introduce una contribuci\'on qu\'imica al comportamiento el\'ectrico de la membrana,
complementando el modelo el\'ectrico puro.

\subsection{Enlace Entre Impedancia y Potencial de Membrana}

El potencial de membrana est\'a directamente relacionado con la acumulaci\'on de carga:

\begin{equation}
Q = C_m V_m
\end{equation}

Dado que la capacitancia de membrana $C_m$ es dependiente de la frecuencia y surge de las
propiedades diel\'ectricas descritas anteriormente, la respuesta de impedancia porta
informaci\'on sobre la din\'amica de carga. As\'i, las mediciones de bioimpedancia proporcionan
acceso indirecto al potencial de membrana a trav\'es de su relaci\'on con el almacenamiento
capacitivo de carga.

\subsection{Contribuci\'on I\'onica Din\'amica}

La membrana no es puramente pasiva. Los canales i\'onicos introducen mecanismos de transporte
activo que modifican la conductividad efectiva. La corriente i\'onica se modela como:
\begin{equation}
I_i = g_i(V_m, t)(V_m - E_i)
\end{equation}

Esto introduce din\'amica no lineal en el sistema, acoplando las propiedades el\'ectricas con
la funci\'on biol\'ogica.

\paragraph{Extensi\'on: ecuaci\'on de Nernst--Planck.}
En una descripci\'on m\'as completa, la conductividad puede
derivarse de ecuaciones de transporte i\'onico
tipo Nernst--Planck:
\begin{equation}
\mathbf{J}_i = -D_i \nabla c_i
- \frac{z_i e D_i}{k_B T}\, c_i \,\nabla \phi
\end{equation}
incorporando difusi\'on, migraci\'on y acoplamiento
electroqu\'imico. Su implementaci\'on completa se deja como
extensi\'on futura, conectando el marco presente con
la literatura de transporte en medios biol\'ogicos.

\subsection{Ecuaci\'on de Membrana Extendida}

La din\'amica de membrana puede describirse por:

\begin{equation}
C_m \frac{dV_m}{dt} +
\sum_i g_i(V_m,t)(V_m - E_i)
= I_{\text{ext}}
\end{equation}

donde la capacitancia $C_m$ y las conductancias $g_i$ se derivan de las propiedades
electrodin\'amicas y diel\'ectricas subyacentes. Esta ecuaci\'on extiende el modelo de impedancia
pasiva hacia un marco fisiol\'ogico dependiente del tiempo.

\subsection{Din\'amica Temporal con Modelo de Hodgkin--Huxley}

La din\'amica completa del potencial de membrana se modela mediante las ecuaciones de
Hodgkin--Huxley. Al incorporar el transporte de sodio y potasio junto con la din\'amica de
canales, el modelo predice fluctuaciones temporales en la impedancia directamente
correlacionadas con cambios en $V_m(t)$, reflejando el estado electrofisiol\'ogico activo
de la c\'elula:

\begin{lstlisting}[language=Python, caption={Implementaci\'on del modelo de Hodgkin--Huxley y c\'alculo de $Z(\omega, t)$.}]
import numpy as np
epsilon0 = 8.854e-12

dt = 0.01;  t = np.arange(0, 50, dt)
C_m = 1.0;  E_Na = 50;  E_K = -77;  E_L = -54.4
g_Na = 120;  g_K = 36;  g_L = 0.3;  I_ext = 10.0

# [C7] Funciones de compuerta con singularidades resueltas (L'Hopital)
def alpha_m(V):
    dV = V+40.0
    if abs(dV)<1e-7: return 1.0
    return 0.1*dV/(1.0-np.exp(-dV/10.0))
def beta_m(V):  return 4*np.exp(-(V+65)/18)
def alpha_h(V): return 0.07*np.exp(-(V+65)/20)
def beta_h(V):  return 1/(1+np.exp(-(V+35)/10))
def alpha_n(V):
    dV = V+55.0
    if abs(dV)<1e-7: return 0.1
    return 0.01*dV/(1.0-np.exp(-dV/10.0))
def beta_n(V):  return 0.125*np.exp(-(V+65)/80)

V = -65*np.ones_like(t)
m = np.zeros_like(t);  h = np.zeros_like(t);  n = np.zeros_like(t)

for i in range(len(t)-1):
    am,bm = alpha_m(V[i]),beta_m(V[i])
    ah,bh = alpha_h(V[i]),beta_h(V[i])
    an,bn = alpha_n(V[i]),beta_n(V[i])
    m[i+1] = m[i] + dt*(am*(1-m[i]) - bm*m[i])
    h[i+1] = h[i] + dt*(ah*(1-h[i]) - bh*h[i])
    n[i+1] = n[i] + dt*(an*(1-n[i]) - bn*n[i])
    I_Na = g_Na*(m[i]**3)*h[i]*(V[i]-E_Na)
    I_K  = g_K*(n[i]**4)*(V[i]-E_K)
    I_L  = g_L*(V[i]-E_L)
    V[i+1] = V[i] + dt*(I_ext - I_Na - I_K - I_L)/C_m

f     = np.logspace(3, 7, 200)
omega = 2*np.pi*f + 1e-20        # [C1]

def cole_cole(eps_inf, eps_s, tau, alpha, omega):
    return eps_inf + (eps_s-eps_inf)/(1+(1j*omega*tau)**(1-alpha))

# Cole-Cole calculado UNA VEZ fuera del loop (eficiencia)
eps_cc = cole_cole(5*epsilon0, 80*epsilon0, 1e-5, 0.2, omega)

def sigma_from_voltage(V_scalar, base=1.0):
    return base + 0.02*(V_scalar + 65)

Z_time = []
for i in range(len(t)):
    sigma    = sigma_from_voltage(V[i])
    # eps_comp ya incluye conduccion: eps_cc - j*sigma/w
    # Z = 1/(jw*eps_comp) evita contar sigma dos veces
    eps_comp = eps_cc - 1j*sigma/omega
    Z_time.append(1/(1j*omega*eps_comp))
# [C8] Sin Z_time.append(Z) duplicado que causaba NameError
Z_time = np.array(Z_time)
\end{lstlisting}

\subsection{Interpretaci\'on Multiscala}

La membrana act\'ua como un puente entre escalas:

\begin{itemize}
\item \textbf{Microsc\'opica}: polarizaci\'on de mol\'eculas y \'atomos
\item \textbf{Mesosc\'opica}: respuesta diel\'ectrica e impedancia
\item \textbf{Macrosc\'opica}: potencial de membrana y corrientes i\'onicas
\end{itemize}

Esta perspectiva multiscala permite la unificaci\'on de las descripciones f\'isica y biol\'ogica
dentro de un \'unico marco te\'orico:

\begin{equation}
\text{Atomic Polarization}
\rightarrow
\alpha_{\text{eff}}
\rightarrow
\varepsilon^*(\omega)
\rightarrow
Z_m(\omega)
\rightarrow
V_m
\rightarrow
I_{\text{ion}}
\rightarrow
\text{Cellular State}
\end{equation}

\subsection{Consideraciones Energ\'eticas}

El sistema de membrana involucra almacenamiento y disipaci\'on de energ\'ia:

\begin{itemize}
\item Energ\'ia capacitiva: $U = \frac{1}{2} C_m V_m^2$
\item Disipaci\'on: corrientes i\'onicas
\item Entrada de energ\'ia: transporte activo (consumo de ATP)

\end{itemize}

\subsection{Marco Multiscala Unificado}

El marco propuesto integra:

\begin{itemize}
\item Polarizaci\'on microsc\'opica (escala at\'omica)
\item Respuesta diel\'ectrica efectiva (escala mesosc\'opica)
\item Impedancia de membrana (interfaz celular)
\item Potencial de membrana (estado electroqu\'imico)
\item Corrientes i\'onicas (din\'amica biol\'ogica)
\end{itemize}

Esto establece una conexi\'on directa entre la f\'isica fundamental y el comportamiento celular,
habilitando el modelado predictivo de sistemas biol\'ogicos.

\subsection{Alcance de la Extensi\'on Fisiol\'ogica}
El presente trabajo no pretende modelar completamente todos los procesos bioqu\'imicos.

En cambio, establece un marco f\'isicamente fundamentado sobre el cual pueden construirse
modelos biol\'ogicos m\'as complejos. Procesos como el bombeo activo de iones, el
comportamiento estoc\'astico de canales y la se\~{n}alizaci\'on intracelular son reconocidos pero
no incluidos
expl\'icitamente. El modelo representa sus efectos colectivos a trav\'es de par\'ametros efectivos

tales como conductancia, tiempos de relajaci\'on y concentraciones i\'onicas. La meta no es
resolver toda la complejidad biol\'ogica, sino representarla en una forma que sea tanto
f\'isicamente significativa como computacionalmente tratable.

\subsection{Implicaciones para el Modelado Predictivo}

Al vincular las propiedades electrodin\'amicas con las variables fisiol\'ogicas, el modelo
abre la posibilidad de predecir estados celulares a partir de mediciones el\'ectricas. Esto
proporciona una base para el desarrollo de herramientas diagn\'osticas no invasivas basadas
en bioimpedancia.

%=================================================================

\section{Pipeline Completo: de la F\'isica al Diagn\'ostico}

\label{sec:pipeline}

%=================================================================

\subsection{Estructura del Sistema}

El pipeline completo del presente trabajo integra:

\begin{enumerate}
\item Simulaci\'on f\'isica (HH + $Z(\omega,t)$)

\item Generaci\'on de base de datos sint\'etica
\item Red neuronal para predicci\'on fisiol\'ogica
\end{enumerate}

\begin{equation}
\text{Simulaci\'on HH} + Z(\omega,t)
\;\rightarrow\;
\text{Dataset sint\'etico}
\;\rightarrow\;
\text{Red neuronal}
\;\rightarrow\;
\text{Predicci\'on fisiol\'ogica}
\end{equation}

\subsection{Observable Macrosc\'opica del Modelo}

La impedancia medible puede expresarse funcionalmente como:

\begin{equation}
Z(\omega) =
\mathcal{F}\left[V_m(t),\,g_i(V_m,t),\,C_m\right]
\end{equation}

Esta relaci\'on establece que la bioimpedancia es una observable multiscala que codifica
din\'amica de cargas microsc\'opica, estructura mesosc\'opica y procesos fisiol\'ogicos activos.

%=================================================================

\section{Resultados e Interpretaci\'on F\'isica}

\label{sec:results}

%=================================================================

\noindent\textit{Esta secci\'on presenta los resultados esperados del modelo y, m\'as importante,
su interpretaci\'on f\'isica profunda. No basta graficar: se requiere conectar cada caracter\'istica
del espectro con un mecanismo f\'isico espec\'ifico, explicar c\'omo cambia $Z$ durante un
potencial de acci\'on, identificar qu\'e frecuencias son m\'as sensibles a $g_{Na}$ y $g_K$, y
se\~{n}alar d\'onde emerge el efecto Maxwell--Wagner.}

\subsection{Magnitud de la Impedancia: Comportamiento en Tres Reg\'imenes}

\begin{figure}[H]
\centering
\framebox[0.78\textwidth]{\rule{0pt}{4.5cm}
\centering $|Z(\omega)|$ vs $f$ (escala log-log) --- marcador de posici\'on}
\caption{Magnitud de impedancia en funci\'on de la frecuencia. Tres reg\'imenes distintos
son visibles: meseta capacitiva a baja $f$, dispersi\'on $\beta$ por Maxwell--Wagner, y
meseta resistiva a alta $f$. Resultados pendientes de validaci\'on experimental.}
\label{fig:Zmag}
\end{figure}

\begin{itemize}

\item \textbf{Baja frecuencia ($f < 10^4\,\text{Hz}$)}: $|Z|$ grande y casi constante.
La corriente fluye completamente por el plasma extracelular; la membrana act\'ua como
circuito abierto. El valor de esta meseta est\'a determinado por $\sigma_3$ y $\phi$.
\item \textbf{Transici\'on ($10^4$--$10^7\,\text{Hz}$)}: $|Z|$ decrece en
aproximadamente un orden de magnitud. La capacitancia de membrana va en cortocircuito
progresivamente. Aqu\'i aparece el efecto Maxwell--Wagner (v\'ease Sec.~\ref{subsec:MW_loc}).
\item \textbf{Alta frecuencia ($f > 10^7\,\text{Hz}$)}: $|Z|$ alcanza una meseta
inferior. Ambos compartimentos conducen en paralelo; la impedancia refleja la
combinaci\'on de $\sigma_1$ y $\sigma_3$.
\end{itemize}

\subsection{Diagrama de Nyquist: Semic\'irculo Deprimido}

\begin{figure}[H]
\centering
\framebox[0.78\textwidth]{\rule{0pt}{4.5cm}
\centering Nyquist: $-\text{Im}(Z)$ vs $\text{Re}(Z)$ --- marcador de posici\'on}
\caption{Diagrama de Nyquist esperado. El semic\'irculo deprimido es la huella
caracter\'istica de la distribuci\'on de tiempos de relajaci\'on celular.
Resultados pendientes.}
\label{fig:nyquist}
\end{figure}

Cada caracter\'istica geom\'etrica del semic\'irculo deprimido codifica un par\'ametro f\'isico
distinto, lo que convierte el diagrama de Nyquist en un mapa de propiedades celulares:

\begin{itemize}

\item \textbf{\'Angulo de depresi\'on} $\theta = \alpha_{CC}\cdot\pi/2$: mide la
heterogeneidad de Cole--Cole (variabilidad en tama\~{n}o y estado celular).
\item \textbf{Frecuencia en el v\'ertice} $f_{MW} = 1/(2\pi\tau_{MW})$: corresponde al
tiempo de relajaci\'on de Maxwell--Wagner; depende del grosor de membrana y la
conductividad del citoplasma.
\item \textbf{Di\'ametro} $R_0 - R_\infty$: proporcional al hematocrito $\phi$; un
descenso de 10\% en $\phi$ (anemia moderada) produce una reducci\'on del di\'ametro
del semic\'irculo.
\item \textbf{Intersecci\'on a baja $f$} ($R_0$): solo conductividad extracelular.
\item \textbf{Intersecci\'on a alta $f$} ($R_\infty$): combinaci\'on paralela de ambos
compartimentos.
\end{itemize}

\subsection{D\'onde Aparece el Efecto Maxwell--Wagner}

\label{subsec:MW_loc}

El efecto Maxwell--Wagner se manifiesta como la transici\'on de $|Z|$ entre las dos mesetas.
F\'isicamente, ocurre cuando la frecuencia angular $\omega$ satisface:

\begin{equation}
\omega \approx \frac{1}{\tau_{MW}} =
\frac{\sigma_{\text{int}} + 2\sigma_3}{\varepsilon_{\text{int}} + 2\varepsilon_3}
\end{equation}

Para los par\'ametros t\'ipicos de eritrocitos ($\sigma_1 = 0.5\,\text{S/m}$, $d = 5\,\text{nm}$), se
estima $f_{MW} \approx 1$--$10\,\text{MHz}$. Por encima de esta
frecuencia, la membrana ya no separa efectivamente los compartimentos y la corriente

penetra al citoplasma. La identificaci\'on de $f_{MW}$ en un espectro medido permite
estimar directamente las propiedades de la membrana sin disociar las c\'elulas.

\subsection{C\'omo Cambia $Z(\omega)$ Durante un Potencial de Acci\'on}

Esta es una de las predicciones m\'as ricas del marco unificado. Durante un potencial de
acci\'on generado por el modelo de Hodgkin--Huxley, ocurre la siguiente secuencia:

\begin{enumerate}
\item \textbf{Fase de reposo} ($V_m \approx -65\,\text{mV}$): $Z(\omega)$ muestra el
semic\'irculo de referencia. La conductividad es baja y la membrana act\'ua como capacitor.
\item \textbf{Despolarizaci\'on r\'apida} ($V_m \uparrow$, canales Na$^+$ abren): el aumento
brusco de $g_{Na}$ incrementa la corriente i\'onica. Esto eleva $\sigma(t)$, lo que
\emph{reduce} $|Z|$ en toda la banda de frecuencia (meseta baja se contrae).
\item \textbf{Pico del potencial} ($V_m \approx +40\,\text{mV}$): m\'axima conductividad
efectiva. El semic\'irculo en el diagrama de Nyquist se aplana y se desplaza hacia la
izquierda.
\item \textbf{Repolarizaci\'on} ($g_K \uparrow$, Na$^+$ inactivan): $\sigma(t)$ decrece
gradualmente. $|Z|$ recupera su valor basal.
\item \textbf{Post-hiperpolarizaci\'on}: el sistema retorna al estado de reposo; $Z(\omega)$
vuelve al semic\'irculo original.
\end{enumerate}

Por tanto, \emph{monitorear $|Z(t)|$ a una frecuencia fija en la regi\'on de dispersi\'on $\beta$
equivale a monitorear indirectamente el estado electrofisiol\'ogico de la c\'elula.}

\subsection{Qu\'e Frecuencias Son M\'as Sensibles a Na$^+$ y K$^+$}

El an\'alisis de sensibilidad revela que los canales i\'onicos afectan preferentemente
diferentes bandas del espectro:

\begin{itemize}
\item \textbf{Canales de Na$^+$ ($g_{Na}$)}: mayor impacto en la \emph{meseta de alta
frecuencia} ($f > 10^6\,\text{Hz}$), donde la corriente penetra el citoplasma. Un
aumento de $g_{Na}$ reduce $|Z|$ en esta banda. La raz\'on f\'isica es que $g_{Na}$
modifica $\sigma(t)$ durante la despolarizaci\'on, que es una propiedad intracelular.
\item \textbf{Canales de K$^+$ ($g_K$)}: afectan principalmente la \emph{fase de
transici\'on} y la repolarizaci\'on. Un $g_K$ mayor produce una repolarizaci\'on m\'as
r\'apida, lo que se traduce en una menor duraci\'on de la excursi\'on de $Z$ en el tiempo.
\item \textbf{Combinaci\'on $g_{Na}/g_K$}: la raz\'on $g_{Na}/g_K$ es el par\'ametro que
m\'as afecta la amplitud del potencial de acci\'on y, por tanto, la amplitud de la
modulaci\'on temporal de $Z(\omega, t)$. Este ratio es recuperable por el modelo inverso.
\end{itemize}

\subsection{Bioimpedancia como Observable Multiscala}

Los resultados confirman que la bioimpedancia es una observable multiscala que no es una
constante intr\'inseca del material, sino una propiedad emergente resultante de procesos
electrodin\'amicos y fisiol\'ogicos acoplados. En particular:

\begin{itemize}
\item A \emph{baja frecuencia}: $Z$ codifica geometr\'ia celular ($\phi$, $a$) y conductividad

del plasma ($\sigma_3$).
\item En la \emph{dispersi\'on $\beta$}: $Z$ codifica las propiedades de membrana ($d$,
$C_m$, $\tau_{MW}$).
\item A \emph{alta frecuencia}: $Z$ codifica conductividad del citoplasma ($\sigma_1$) y
el estado de los canales i\'onicos ($g_{Na}$, $g_K$).
\item En el \emph{dominio temporal}: $Z(t)$ codifica la excitabilidad celular y la
din\'amica del potencial de acci\'on.
\end{itemize}

\subsection{Sensibilidad a Par\'ametros Fisiol\'ogicos}

\begin{table}[H]
\centering
\caption{Sensibilidad del espectro de impedancia a cada par\'ametro f\'isico con
identificaci\'on de la banda espectral y mecanismo f\'isico responsable.}
\label{tab:sensibilidad}
\begin{tabular}{@{}llll@{}}
\toprule
\textbf{Par\'ametro} & \textbf{Sensibilidad} & \textbf{Banda espectral} & \textbf{Mecanismo} \\
\midrule
$\phi$ (hematocrito) & Alta & Espectro completo & Reducci\'on de caminos de corriente \\
$\sigma_3$ (plasma) & Alta & $f < 10^4\,\text{Hz}$ & Conducci\'on extracelular \\
$d$ (membrana) & Alta & Entorno de $f_{MW}$ & Capacitancia de membrana \\
$\sigma_1$ (citoplasma) & Alta & $f > 10^6\,\text{Hz}$ & Conducci\'on intracelular \\
$g_{Na}$ & Media & $f > 10^5\,\text{Hz}$ & Conductividad din\'amica Na$^+$ \\
$g_K$ & Media & Tiempo: repolarizaci\'on & Duraci\'on del potencial de acci\'on \\
$\alpha_{CC}$ & Media & Ancho de la transici\'on & Distribuci\'on de $\tau_{MW}$ \\

$a$ (radio celular) & Media & Entorno de $f_{MW}$ & Cambia $\tau_{MW} \propto a$ \\
$\varepsilon_1$ & Baja & $f > 10^8\,\text{Hz}$ & Polarizaci\'on electr\'onica \\
$\varepsilon_2$ & Baja & Cerca de $f_{MW}$ & Almacenamiento capacitivo \\
\bottomrule
\end{tabular}
\end{table}

\subsection{Efecto de la Heterogeneidad Celular}

Introducir una distribuci\'on gaussiana de tama\~{n}os celulares conduce a un ensanchamiento
del espectro de impedancia, consistente con el comportamiento de Cole--Cole ($\alpha_{CC} > 0$).
El par\'ametro $\alpha_{CC}$ actua como un biomarcador de heterogeneidad: poblaciones
celulares m\'as heterog\'eneas (por ejemplo, sangre con anemia) presentan valores mayores de
$\alpha_{CC}$.

La variabilidad biol\'ogica no es ruido; es informaci\'on sobre el estado de la poblaci\'on
celular.

\subsection{Validaci\'on de L\'imites Anal\'iticos Conocidos}
El modelo es internamente consistente con los l\'imites anal\'iticos:

\begin{itemize}
\item $\phi \to 0$: $Z(\omega) \to 1/\sigma_3$ (medio homog\'eneo de plasma puro).
$\checkmark$
\item $d \to \infty$ (membrana gruesa): la c\'elula actua como esfera aislante a toda

frecuencia. $\checkmark$
\item $\sigma_2 \to \sigma_1$ (membrana conductora): desaparece la dispersi\'on $\beta$.
$\checkmark$
\item $\omega \to 0$: $Z \to R_0 = 1/[\sigma_3(1 - 3\phi/2)]$ (Maxwell--Garnett est\'atico).
$\checkmark$
\end{itemize}

\subsection{Hacia un Modelo Fisiol\'ogico M\'as Realista}

Para abordar las limitaciones fisiol\'ogicas, el modelo puede extenderse para incluir:

\begin{itemize}
\item Propiedades heterog\'eneas de membrana (variabilidad local de $C_m$)
\item Modulaci\'on din\'amica de canales i\'onicos dependiente de voltaje
\item Tiempos de relajaci\'on distribuidos con espectro continuo
\item Osmolaridad real y efecto de la difusi\'on (Nernst--Planck)
\end{itemize}

Estas extensiones mejoran el realismo biol\'ogico mientras preservan la interpretabilidad
f\'isica. Toda la complejidad adicional se introduce a trav\'es de par\'ametros medibles,
garantizando que el modelo permanezca experimentalmente testeable.

%=================================================================

\section{Estrategia de Validaci\'on Experimental}

\label{sec:validacion}

%=================================================================

\noindent\textit{El modelo f\'isico y el pipeline de aprendizaje autom\'atico se han desarrollado
sobre una base \'unicamente de simulaci\'on. Esta secci\'on describe el protocolo experimental
concreto planificado para validar tanto el modelo directo como el inverso.}

\subsection{Protocolo de Medici\'on}

\textbf{Preparaci\'on de muestras.} Se obtendr\'a sangre humana completa de donantes sanos
bajo aprobaci\'on \'etica institucional. Las muestras se preparar\'an a hematocritos controlados
$\phi \in \{0.20, 0.30, 0.40, 0.45, 0.50\}$ por centrifugaci\'on y resuspensi\'on en plasma
autól\'ogo. La conductividad del plasma $\sigma_3$ se medir\'a de forma independiente.

\textbf{Sistema de medici\'on.} Se usar\'a una configuraci\'on de cuatro electrodos (tetrapolar)
para eliminar los artefactos de polarizaci\'on de electrodo \citep{Sanchez2013}. Las
mediciones se realizar\'an sobre $f \in [10^3, 10^8]\,\text{Hz}$ (100 puntos equiespaciados
logar\'itmicamente), cubriendo la dispersi\'on $\beta$ completa.

\subsection{Validaci\'on del Modelo}

El modelo directo se validar\'a comparando $|Z(\omega)|$ simulado y el diagrama de Nyquist
contra espectros medidos a valores conocidos de $(\phi, \sigma_3)$. La figura de m\'erito es
el error cuadr\'atico medio normalizado (NRMSE):

\begin{equation}
\text{NRMSE} = \frac{1}{N_f}\sqrt{\sum_{k=1}^{N_f}
\left(\frac{|Z_{\text{sim}}(\omega_k)| - |Z_{\text{med}}(\omega_k)|}
{|Z_{\text{med}}(\omega_k)|}\right)^2}
\end{equation}

Se considera aceptable un objetivo de NRMSE $< 5\%$ en la banda de dispersi\'on $\beta$.

\begin{table}[H]
\centering
\caption{Protocolo de validaci\'on: condiciones experimentales, m\'etodos de referencia
y criterios de aceptaci\'on.}
\label{tab:validacion}
\begin{tabular}{@{}llll@{}}
\toprule
\textbf{Objetivo} & \textbf{M\'etodo de referencia} & \textbf{Criterio}
& \textbf{N muestras} \\
\midrule
Modelo directo & $|Z|$ medido vs.\ simulado & NRMSE $< 5\%$ & 5 niveles $\phi$ \\
$\hat{\phi}$ & Microhematocrito centrifugado & MAPE $< 8\%$ & $\geq 30$ \\
$\hat{\sigma}_1$ & Hem\'olisis + conduct\'imetro & MAPE $< 10\%$ & $\geq 20$ \\
$\hat{\sigma}_3$ & Medici\'on directa & MAPE $< 8\%$ & $\geq 30$ \\
\bottomrule

\end{tabular}
\end{table}

Adicionalmente, la consistencia del modelo se validar\'a a trav\'es de:

\begin{itemize}
\item Comportamiento en l\'imites anal\'iticos conocidos (baja frecuencia $\to$ capacitivo,
alta frecuencia $\to$ resistivo)
\item Estabilidad bajo variaciones de par\'ametros
\item Consistencia con datos publicados \citep{Schwan1957, Asami2002, Buendia2015}
\end{itemize}

%=================================================================

\section{Limitaciones del Modelo}


\label{sec:limitations}

%=================================================================
\paragraph{Alcance declarado del modelo.}
El modelo no pretende describir \emph{todos} los procesos
fisiol\'ogicos, sino capturar los \textbf{mecanismos
electromagn\'eticos dominantes} relevantes para la
bioimpedancia en $10^2$--$10^9$\,Hz.
No incluye interacciones moleculares detalladas ni
procesos bioqu\'imicos intracelulares, los cuales se
consideran fuera del alcance del presente trabajo.

El modelo de aprendizaje autom\'atico no reemplaza la
f\'isica del sistema: act\'ua como herramienta para el
problema inverso, cuya generalizaci\'on a datos reales
depender\'a de la fidelidad del modelo f\'isico
(l\'inea prioritaria de trabajo futuro).



Una evaluaci\'on rigurosa de las limitaciones del modelo es esencial para su aplicaci\'on
responsable.
\paragraph{Geometr\'ia esf\'erica.}

Los eritrocitos son discos bic\'oncavos ($\sim 8\,\mu$m de di\'ametro, $\sim 2\,\mu$m de
grosor). La aproximaci\'on esf\'erica preserva la f\'isica electrost\'atica esencial pero no puede
reproducir la polarizabilidad dependiente de la orientaci\'on de c\'elulas en forma de disco.

\paragraph{Teor\'ia del medio efectivo para suspensiones diluidas.}

La teor\'ia de Maxwell--Garnett es rigurosamente v\'alida solo para $\phi \ll 1$. A hematocritos
fisiol\'ogicos ($\phi \approx 0.40$), las interacciones c\'elula-c\'elula se vuelven significativas.

\paragraph{Citoplasma y membrana homog\'eneos.}

El citoplasma se modela como un medio uniforme con una sola relajaci\'on de Cole--Cole,
ignorando el n\'ucleo, las mitocondrias y los filamentos del citoesqueleto.

\paragraph{R\'egimen de respuesta lineal.}

Todas las ecuaciones asumen respuesta diel\'ectrica lineal, v\'alida para los campos de medici\'on
t\'ipicos de bioimpedancia ($\sim 10\,\text{mV/cm}$). Se rompe cerca de los umbrales de
electroporaci\'on ($\sim 1\,\text{kV/cm}$).

\paragraph{Linealidad de la din\'amica de canales.}

El presente modelo se enfoca en las propiedades el\'ectricas electrodin\'amicas y pasivas de la
membrana, incorporando corrientes i\'onicas simplificadas. Procesos activos como el bombeo de
iones, el comportamiento estoc\'astico de canales y la regulaci\'on bioqu\'imica intracelular no
se modelan expl\'icitamente, pero pueden incorporarse en extensiones futuras.

\paragraph{Limitaciones biol\'ogicas.}

El modelo asume c\'elulas esf\'ericas y medios homog\'eneos, lo que puede limitar su aplicabilidad
a tejidos biol\'ogicos complejos con anisotrop\'ia, interacciones c\'elula-c\'elula o din\'amica de
fluidos. Extender el modelo a geometr\'ias no esf\'ericas y tejidos heterog\'eneos es un objetivo
prioritario del trabajo futuro.
\paragraph{Generalizaci\'on del aprendizaje autom\'atico.}

La red neuronal se entrena \'unicamente con datos de simulaci\'on. La transferencia a datos
experimentales requiere que los artefactos instrumentales (polarizaci\'on de electrodo,
capacitancia par\'asita) sean corregidos en el preprocesamiento.

%=================================================================

\section{Cadena Conceptual Completa}

\label{sec:chain}

%=================================================================

El marco multiscala completo se resume en la siguiente cadena:

\begin{equation}
\underbrace{\text{MC}}_{\S\ref{sec:atomic}}
\;\xrightarrow{\text{Cls.-Mos.}}\;
\underbrace{\alpha(\omega)}_{\S\ref{sec:atomic}}
\;\xrightarrow{\text{Cole--Cole}}\;
\underbrace{\varepsilon^*(\omega)}_{\S\ref{sec:debye}}
\;\xrightarrow{\text{MW}}\;
\underbrace{Z_{\text{c\'el.}}}_{\S\ref{sec:cell}}
\;\xrightarrow{\text{MG}}\;
\underbrace{Z_{\text{sangre}}}_{\S\ref{sec:sangre}}
\;\xrightarrow{\text{MC}}\;
\underbrace{\text{BD}}_{\S\ref{sec:ml}}
\;\xrightarrow{\text{RN}}\;
\underbrace{\text{Predicci\'on}}_{\S\ref{sec:results}}
\end{equation}

La cadena extendida que incorpora las extensiones fisiol\'ogicas es:

\begin{equation}
\alpha_{\text{eff}}
\rightarrow
\varepsilon^*(\omega)
\rightarrow
Z_m(\omega)
\rightarrow
V_m
\rightarrow
I_{\text{ion}}
\rightarrow
\text{Cellular State}
\end{equation}

Esta cadena destaca la naturaleza multiscala del problema: cada nivel proporciona la entrada
necesaria y suficiente para el siguiente. La consistencia f\'isica se preserva en todo momento
--- no se introducen par\'ametros fenomenol\'ogicos sin derivaci\'on.

%=================================================================

\section{Aplicaciones}

\label{sec:aplicaciones}

%=================================================================

\begin{itemize}
\item \textbf{An\'alisis de sangre}: estimaci\'on no invasiva del hematocrito y la
conductividad del citoplasma desde un sensor en la mu\~{n}eca o en el dedo.
\item \textbf{Detecci\'on de enfermedades}: los cambios en la capacitancia de la membrana
y la conductividad del citoplasma son biomarcadores conocidos para la anemia,
la malaria, la anemia falciforme y ciertos tipos de c\'ancer.
\item \textbf{Caracterizaci\'on de tejidos}: el marco se aplica directamente a cualquier
tejido modelado como una suspensi\'on de c\'elulas en matriz extracelular.
\item \textbf{Microflu\'idica de c\'elula \'unica}: en canales microflu\'idicos, el modelo puede
adaptarse para predecir propiedades individuales de c\'elulas a partir de impulsos
de impedancia de tr\'ansito.
\item \textbf{Monitoreo de excitabilidad celular}: la bioimpedancia dependiente del tiempo,
vinculada con la din\'amica de Hodgkin--Huxley, puede usarse para monitorear
estados de excitabilidad celular de forma no invasiva.
\end{itemize}

%=================================================================

\section{Aportaciones Principales}

\label{sec:contributions}

%=================================================================

Las principales contribuciones de este trabajo pueden resumirse como sigue.

La contribuci\'on principal de este trabajo es la integraci\'on de un modelo
multiescala que conecta la f\'isica fundamental con la fisiolog\'ia celular,
permitiendo la construcci\'on de un marco predictivo basado en bioimpedancia.
A diferencia de modelos tradicionales que emplean descripciones fenomenol\'ogicas,
el enfoque propuesto integra diferentes niveles de descripci\'on ---desde la
polarizaci\'on at\'omica hasta la din\'amica de membranas celulares--- incorporando
adem\'as efectos de distribuci\'on y heterogeneidad biol\'ogica.

El aporte central radica en el acoplamiento consistente entre modelos f\'isicos de polarizaci\'on
at\'omica y modelos fisiol\'ogicos de transporte i\'onico, lo cual permite una descripci\'on
integrada del sistema desde la escala at\'omica hasta la respuesta el\'ectrica macrosc\'opica. A
diferencia de los enfoques fenomenol\'ogicos existentes, cada par\'ametro del modelo posee un
significado f\'isico preciso y puede relacionarse directamente con una magnitud medible.

\begin{itemize}

\item Desarrollo de un marco te\'orico multiescala que conecta la polarizaci\'on at\'omica,
la respuesta diel\'ectrica y la bioimpedancia, derivando cada nivel desde primeros principios.

\item Derivaci\'on de una formulaci\'on din\'amica de bioimpedancia $Z(\omega,t)$ que acopla
la electrodin\'amica dependiente de la frecuencia con la din\'amica temporal del potencial de
membrana.

\item Integraci\'on de la din\'amica de canales i\'onicos de Hodgkin--Huxley en modelos
diel\'ectricos y conductivos de medios biol\'ogicos, unificando biolog\'ia celular y f\'isica
del continuo.

\item Construcci\'on de un dataset enteramente f\'isicamente fundamentado para entrenar
modelos
de aprendizaje autom\'atico, evitando la generaci\'on emp\'irica o ad-hoc de datos.
\item Demostraci\'on de que los par\'ametros fisiol\'ogicos pueden inferirse a partir de
espectros de impedancia usando redes neuronales, validando la invertibilidad del modelo.

\item Establecimiento de un pipeline computacional que vincula teor\'ia, simulaci\'on y
predicci\'on basada en datos dentro de un marco unificado.

\end{itemize}

%=================================================================

\section{Archivos Generados por la Implementaci\'on}

\label{sec:archivos}

%=================================================================

La ejecuci\'on completa de los cuatro scripts produce los siguientes archivos,
que sirven tanto como resultados del trabajo como insumos para futuros an\'alisis:

\begin{center}
\renewcommand{\arraystretch}{1.35}
\begin{tabular}{@{}llp{7cm}@{}}
\toprule
\textbf{Archivo} & \textbf{Tipo} & \textbf{Contenido} \\
\midrule
\texttt{bioimpedancia.db}            & SQLite & 5\,000 espectros etiquetados con par\'ametros fisiol\'ogicos \\
\texttt{bionet\_inverso.pt}          & PyTorch & Pesos de BioNet + estad\'isticas de normalizaci\'on \\
\texttt{y\_stats.npy}                & NumPy & Media y desviaci\'on est\'andar de etiquetas (para desnormalizar) \\
\texttt{01\_espectros\_simulacion.pdf} & Figura & $|Z|$, fase y diagrama de Nyquist para par\'ametros base \\
\texttt{01\_barrido\_parametros.pdf}   & Figura & Sensibilidad de $|Z(\omega)|$ a $\sigma_1$, $\sigma_3$, $\phi$ \\
\texttt{03\_curva\_aprendizaje.pdf}    & Figura & MSE de entrenamiento y validaci\'on por \'epoca \\
\texttt{04a\_predicho\_vs\_real.pdf}   & Figura & Dispersi\'on $\hat{y}$ vs $y$ para los 5 par\'ametros \\
\texttt{04b\_histograma\_errores.pdf}  & Figura & Distribuci\'on del error relativo por par\'ametro \\
\texttt{04c\_mapa\_correlacion.pdf}    & Figura & Mapa de calor $|r|$ entre features y par\'ametros \\
\texttt{04e\_espectros\_parametros.pdf}& Figura & Barrido de $|Z|$ en rangos fisiol\'ogicos \\
\bottomrule
\end{tabular}
\end{center}

%=================================================================

\section{Conclusiones}

\label{sec:conclusions}

%=================================================================

En este trabajo se ha desarrollado un modelo multiescala que permite conectar
la polarizaci\'on microsc\'opica de la materia con la respuesta macrosc\'opica de
bioimpedancia en sistemas biol\'ogicos. A trav\'es de una construcci\'on progresiva,
se integraron conceptos de electrodim\'amica, mec\'anica estad\'istica y fisiolog\'ia
celular, dando lugar a un marco te\'orico consistente y f\'isicamente fundamentado.

Partiendo de la descripci\'on de la interacci\'on campo-materia a nivel at\'omico,

se estableci\'o la polarizabilidad como un elemento clave para comprender la
respuesta diel\'ectrica. Mediante la relaci\'on de Clausius--Mossotti, se logr\'o
conectar esta descripci\'on microsc\'opica con propiedades macrosc\'opicas del medio,
permitiendo introducir posteriormente efectos din\'amicos a trav\'es de modelos
de relajaci\'on como Debye y Cole--Cole.

La extensi\'on del modelo hacia sistemas biol\'ogicos incorpor\'o la estructura y
din\'amica de la membrana celular, as\'i como fen\'omenos de polarizaci\'on interfacial
descritos por el mecanismo de Maxwell--Wagner. De manera importante, se integr\'o
la contribuci\'on de canales i\'onicos, lo cual permiti\'o capturar efectos
fisiol\'ogicos relevantes y establecer una dependencia din\'amica en la conductividad.

Uno de los resultados m\'as relevantes de este trabajo es la interpretaci\'on de la
bioimpedancia como un observable f\'isico que contiene informaci\'on sobre la
estructura y din\'amica interna del sistema. Esta perspectiva permite reformular
el problema como uno de tipo inverso, en el cual se busca inferir par\'ametros
fisiol\'ogicos a partir de mediciones el\'ectricas.

Debido a la complejidad y no linealidad de este problema, se implement\'o un enfoque
basado en aprendizaje autom\'atico, el cual demostr\'o ser capaz de aproximar la
relaci\'on entre la respuesta en frecuencia y los par\'ametros internos del sistema.
Esto abre la posibilidad de desarrollar herramientas predictivas para la
caracterizaci\'on no invasiva de tejidos biol\'ogicos.

Si bien el modelo propuesto se basa en ciertas idealizaciones, como geometr\'ias
simplificadas y medios homog\'eneos, proporciona una base s\'olida para futuras
extensiones que incluyan mayor complejidad estructural y validaci\'on experimental.

En conjunto, este trabajo establece un puente entre la f\'isica fundamental y la

fisiolog\'ia, proponiendo una nueva forma de interpretar la bioimpedancia como
una herramienta capaz de revelar informaci\'on interna de sistemas biol\'ogicos.
Este enfoque no solo contribuye a una mejor comprensi\'on de los fen\'omenos
involucrados, sino que tambi\'en abre nuevas oportunidades en el desarrollo de
t\'ecnicas diagn\'osticas basadas en mediciones el\'ectricas.

Como trabajo futuro, se propone la incorporaci\'on de geometr\'ias m\'as realistas,
modelos m\'as detallados de transporte i\'onico y la validaci\'on experimental con
datos reales, lo cual permitir\'a evaluar el alcance y aplicabilidad del modelo
en contextos biom\'edicos.

\paragraph{Fundamento f\'isico.}
La polarizabilidad at\'omica se deriv\'o desde primeros principios usando el
oscilador de Lorentz y la teor\'ia de perturbaciones cu\'antica. El par\'ametro
$\tau$ no es \'unicamente una constante matem\'atica, sino que representa procesos
f\'isicos como la rotaci\'on de dipolos o la reorganizaci\'on estructural del medio.
De manera similar, la conductividad $\sigma$ est\'a asociada al movimiento de iones,
particularmente sodio y potasio, conectando el modelo directamente con la
fisiolog\'ia del sistema.

\paragraph{Relajaci\'on din\'amica y heterogeneidad.}
La generalizaci\'on de Cole--Cole captura la relajaci\'on distribuida mediante
$\alpha_{CC}$, un posible biomarcador de la heterogeneidad biol\'ogica del tejido.
La introducci\'on de una distribuci\'on de tama\~{n}os celulares permite capturar
la variabilidad biol\'ogica inherente a sistemas reales, siendo responsable del
ensanchamiento observado en la respuesta en frecuencia.

\paragraph{Electrodin\'amica unicelular y sangre.}
La soluci\'on cuasiestática de Maxwell identifica $\tau_{MW}$ como el par\'ametro clave

que vincula las propiedades de la membrana con la dispersi\'on $\beta$. La
teor\'ia de Maxwell--Garnett captura la respuesta colectiva de la sangre en
funci\'on del hematocrito, la capacitancia de la membrana y la conductividad del plasma.

\paragraph{Canales i\'onicos y potencial de membrana.}
Los canales i\'onicos introducen una dependencia temporal y no lineal en $\sigma$,
ya que su apertura y cierre dependen de $V_m$. A su vez, el potencial de
membrana modifica la distribuci\'on de cargas en la interfaz celular, afectando
directamente la polarizaci\'on interfacial y, por consiguiente, la impedancia medida.
Esto confirma que cambios fisiol\'ogicos pueden detectarse indirectamente mediante
bioimpedancia.

\paragraph{Aprendizaje autom\'atico como extensi\'on f\'isica.}
La red neuronal no es una caja negra, sino el inverso matem\'atico del modelo
f\'isico. Su capacidad de generalizaci\'on est\'a garantizada por el fundamento
f\'isico de los datos de entrenamiento.

\paragraph{Perspectiva.}
Los pr\'oximos pasos son: (i)~validaci\'on experimental con eritrocitos humanos;
(ii)~geometr\'ias celulares no esf\'ericas; (iii)~polarizaci\'on de electrodo;
(iv)~implementaci\'on en tiempo real para diagn\'ostico en punto de atenci\'on;
y (v)~transferencia del modelo a datos experimentales reales.

% ---- CIERRE FINAL ---
\begin{center}
\textit{
Este trabajo establece un puente entre la f\'isica fundamental y la fisiolog\'ia,
demostrando que la bioimpedancia puede ser utilizada no solo como herramienta
de medici\'on, sino como un medio para acceder a la din\'amica interna de sistemas biol\'ogicos, abriendo nuevas oportunidades en el desarrollo de t\'ecnicas
diagn\'osticas basadas en mediciones el\'ectricas.
}
\end{center}

\section{Impacto y Perspectivas Futuras}

\label{sec:impact}

%=================================================================

El marco desarrollado en este trabajo tiene implicaciones potenciales en m\'ultiples \'areas
de la ciencia y la ingenier\'ia biom\'edica.

En diagn\'ostico, abre la posibilidad de extraer informaci\'on a nivel celular a partir de
mediciones de impedancia no invasivas, con aplicaciones en an\'alisis de sangre, detecci\'on
de enfermedades y caracterizaci\'on de tejidos.

En modelado computacional, proporciona un modelo para integrar leyes f\'isicas con
aprendizaje autom\'atico, garantizando interpretabilidad y consistencia en enfoques
basados en datos.

El trabajo futuro puede extender este modelo incorporando:
\begin{itemize}
\item Geometr\'ias celulares complejas y membranas anisotr\'opicas
\item Interacciones multicelulares y organizaci\'on a nivel de tejido
\item Acoplamiento con din\'amica de fluidos en sistemas capilares
\item Validaci\'on experimental usando mediciones reales de bioimpedancia
\end{itemize}

Estas direcciones fortalecer\'an a\'un m\'as la conexi\'on entre los modelos te\'oricos y las
aplicaciones biom\'edicas del mundo real.

%=================================================================

\section{Trabajo Futuro}

\label{sec:futuro}

%=================================================================

\begin{itemize}
\item Compartimentos subcelulares (n\'ucleo, mitocondrias) como capas conc\'entricas
adicionales en el modelo celular.
\item Geometr\'ias no esf\'ericas: esferoides prolatos/oblatos para c\'elulas nucleadas y
discos bic\'oncavos para eritrocitos.

\item Teor\'ia de Bruggeman o de medio efectivo diferencial para suspensiones
concentradas ($\phi > 0.5$).
\item Acoplamiento con la din\'amica de flujo microflu\'idico y capilar.
\item Incorporaci\'on del modelo de Hodgkin--Huxley completo con compuertas de canal
$m$, $h$, $n$ derivadas desde primeros principios.
\item Modelo de bomba Na$^+$/K$^+$ para sistemas de no-equilibrio.
\item Transferencia de aprendizaje de datos de simulaci\'on a datos experimentales.
\item Validaci\'on experimental: mediciones de bioimpedancia de cuatro electrodos en
sangre humana a hematocritos controlados.
\end{itemize}
%=================================================================

\begin{thebibliography}{99}
\bibitem[Fricke(1925)]{Fricke1925}
Fricke, H.\ (1925).
La capacidad el\'ectrica de suspensiones con referencia especial a la sangre.
\textit{Journal of General Physiology}, \textbf{9}, 137--152.

\bibitem[Schwan(1957)]{Schwan1957}
Schwan, H.\ P.\ (1957).
Propiedades el\'ectricas de tejidos y suspensiones celulares.
\textit{Advances in Biological and Medical Physics}, \textbf{5}, 147--209.

\bibitem[Schwan(1994)]{Schwan1994}
Schwan, H.\ P.\ (1994).
Propiedades el\'ectricas de tejidos y suspensiones celulares: mecanismos y modelos.
\textit{Proc.\ 16th Annual International Conference of the IEEE EMBS}, A70--A71.

\bibitem[Cole \& Cole(1940)]{Cole1940}

Cole, K.\ S.\ y Cole, R.\ H.\ (1940).
Dispersi\'on y absorci\'on en diel\'ectricos {I}: Caracter\'isticas de corriente alterna.
\textit{Journal of Chemical Physics}, \textbf{9}, 341--351.

\bibitem[Foster \& Schwan(1989)]{Foster1989}
Foster, K.\ R.\ y Schwan, H.\ P.\ (1989).
Propiedades diel\'ectricas de tejidos y materiales biol\'ogicos: una revisi\'on cr\'itica.
\textit{Critical Reviews in Biomedical Engineering}, \textbf{17}, 25--104.

\bibitem[Asami(2002)]{Asami2002}
Asami, K.\ (2002).
Caracterizaci\'on de c\'elulas biol\'ogicas por espectroscop\'ia diel\'ectrica.
\textit{Journal of Non-Crystalline Solids}, \textbf{305}, 268--277.

\bibitem[Kaatze(1989)]{Kaatze1989}
Kaatze, U.\ (1989).
Permitividad compleja del agua en funci\'on de la frecuencia y la temperatura.
\textit{Journal of Chemical and Engineering Data}, \textbf{34}, 371--374.

\bibitem[Maxwell(1873)]{Maxwell1873}
Maxwell, J.\ C.\ (1873).
\textit{Tratado sobre Electricidad y Magnetismo}.
Clarendon Press, Oxford.

\bibitem[Wagner(1914)]{Wagner1914}
Wagner, K.\ W.\ (1914).
Explicaci\'on de los procesos de pospolarizaci\'on diel\'ectrica seg\'un las ideas de Maxwell.
\textit{Archiv f\"{u}r Elektrotechnik}, \textbf{2}, 371--387.

\bibitem[Griffiths(1999)]{Griffiths1999}
Griffiths, D.\ J.\ (1999).
\textit{Introducci\'on a la Electrodin\'amica} (3.\textordfeminine~ed.).
Prentice Hall.

\bibitem[S\'anchez et al.(2013)]{Sanchez2013}
S\'anchez, B., Vandersteen, G., Bragos, R.\ y Schoukens, J.\ (2013).
Fundamentos de mediciones de espectroscop\'ia de impedancia de banda ancha usando
excitaciones peri\'odicas.
\textit{Measurement Science and Technology}, \textbf{23}, 105501.
\bibitem[Buend\'ia et al.(2015)]{Buendia2015}
Buend\'ia, R., S\'anchez, B., Gil-Pita, R.\ y Seoane, F.\ (2015).
Validaci\'on experimental del modelo de impedancia de capa \'unica de eritrocitos usando
mediciones en sangre humana.
\textit{Journal of Physics: Conference Series}, \textbf{224}, 012128.

\bibitem[Barsoukov \& Macdonald(2005)]{Barsoukov2005}
Barsoukov, E.\ y Macdonald, J.\ R.\ (eds.) (2005).
\textit{Espectroscop\'ia de Impedancia: Teor\'ia, Experimento y Aplicaciones}
(2.\textordfeminine~ed.).
Wiley-Interscience.

\bibitem[Hanai(1960)]{Hanai1960}
Hanai, T.\ (1960).
Teor\'ia de la dispersi\'on diel\'ectrica debida a la polarizaci\'on interfacial y su
aplicaci\'on a emulsiones.
\textit{Kolloid-Zeitschrift}, \textbf{171}, 23--31.

\bibitem[Pauly \& Schwan(1959)]{Pauly1959}

Pauly, H.\ y Schwan, H.\ P.\ (1959).
Sobre la impedancia de una suspensi\'on de part\'iculas esf\'ericas con una capa.
\textit{Zeitschrift f\"{u}r Naturforschung B}, \textbf{14}, 125--131.

\bibitem[Hodgkin \& Huxley(1952)]{HodgkinHuxley1952}
Hodgkin, A.\ L.\ y Huxley, A.\ F.\ (1952).
Una descripci\'on cuantitativa de la corriente de membrana y su aplicaci\'on a la conducci\'on
y excitaci\'on en el nervio.
\textit{Journal of Physiology}, \textbf{117}, 500--544.
\bibitem[Goldman(1943)]{Goldman1943}
Goldman, D.\ E.\ (1943).
Potencial, impedancia y rectificaci\'on en membranas.
\textit{Journal of General Physiology}, \textbf{27}, 37--60.
\end{thebibliography}

%=================================================================

\end{document}
